% Tabelle der offenen Positionen am Desktop (Reiter Active, blaues Design) – Vektor-Nachbau
% VOLLSTÄNDIGE Fassung ohne Bruchkante, für eine Querformat-Seite (Satzbreite ca. 257 mm).
% Gegenstück zum Ausschnitt positionen-desktop.tex (gleiche Technik, Farben, Platzhalter).
% Grundlage: DOM-Messung originals/messungen/positionen-desktop.json (CSS-px, Fenster
% 2040 x 986; anonymisiert: Ziffern = 0, Instrumente = XXXXXX). Koordinaten unten in
% CSS-px, y nach unten, Ursprung = linke obere Ecke des Blocks.
%
% Gemessen (Original-x):
%   .maintable   1704 x 271, Fläche rgb(4,39,60), Radius 7 px, Inter 400, 12 px / 14 px
%   Leiste oben links (Lage im Block, unabhängig von den Spalten):
%     a.button.bordered  „New order“ x10 y10,2 100 x 29,6, Rand 0,8 px rgb(83,188,81),
%                        Radius 5 px, Fläche frei, padding 6px 10px, Inter 600 weiß,
%                        Zeile 16 px, text-align center
%       i.icon-plus      x22,3 y17,8 12 x 12, icomoon 12 px, U+E91A, weiß
%     input.form-control x120 y10 145 x 30, Fläche und Rand rgb(0,19,40), Radius 4 px,
%                        padding 6px 30px 6px 10px, Zeile 16 px; Platzhalter „Search“
%       i.icon-search    x238 y17,5 15 x 15, icomoon 15 px, U+E915, rgb(168,187,190)
%     ng-select-container x275 / x430 y10 145 x 30, Fläche und Rand rgb(0,19,40),
%                        Radius 4 px; ng-value x275,8 / x430,8, padding 6px 10px,
%                        Zeile 18 px, weiß: „All sides“ (67,4), „All profits“ (73,8)
%   ul.nav-tabs  x1436,3 y10 257,7 x 30, Fläche rgb(0,19,40), Radius 3 px; endet bündig
%                mit der Tabelle (x1694 = 1704 - 10)
%     a.nav-link  flex, padding 0 16px, Text mittig, weiß, capitalize:
%                „Active (0)“ 87 (aktiv: Fläche rgb(14,80,114), Radius 3 0 0 3),
%                „Pending (0)“ 98,1, „History“ 72,6 (Radius 0 3 3 0)
%   th           y47, 28,8 hoch, padding 4px 10px, Text rgb(168,187,190), capitalize;
%                span (inline-block, 20 px Zeile) mit padding-right 15 px; Text beginnt am
%                gemessenen span-x: ID 102,5, Instrument 177, Type 298, Amount 383,8,
%                Trade volume 493,5, Open Rate 640, Open Time 787,5, S/L 941, T/P 1030,9,
%                Swap 1113,6, Commission 1207,4, Total Profit 1344,2, Margin 1467,9
%   td           Zeilen y75,8 / y105,4, je 29,6 hoch, border-top 1,6 px rgb(4,39,60),
%                Fläche rgba(14,80,114,.1), Text mittig (padding 4px 10px); erste Zelle
%                Radius 3 0 0 3, letzte 0 3 3 0
%                Spalten: (Bild 10 w70) | ID 80 w71,8 | Instrument 151,8 w126,2 |
%                Type 278 w83,2 | Amount 361,2 w104,3 | Trade volume 465,5 w149,5 |
%                Open Rate 615 w124,2 | Open Time 739,2 w173,8 | S/L 912,9 w90 |
%                T/P 1002,9 w90 | Swap 1093 w87,8 | Commission 1180,8 w138,1 |
%                Total Profit 1318,9 w127,2 | Margin 1446 w98 | (Zahnrad 1544 w50) |
%                (Schließen 1594 w100)
%   Total Profit grün rgb(83,188,81) (color-green-main) bzw. rot rgb(249,64,86)
%                (color-red-main)
% Grundlinien (Chrome, Gerätepixel 0,8 px: Ascent 12, Descent 3,2 px; icomoon 12 px:
% Ascent 11,2, Descent 0,8 px): Zeilenoberkante + (Zeilenhöhe - 15,2) / 2 + 12.
%   Kopf 65,4; Zeilen 95,0 und 124,6 (belegt durch €-span und Bildmitte, siehe
%   positionen-desktop.tex); Reiter 29,2 (mittig angenommen, nicht gemessen);
%   New order 29,0 (belegt: Oberkante icon-plus 17,8 + Ascent 11,2; der halbe Durchschuss
%   0,4 px fällt weg); Search 10 + 0,8 + 6 + 0,4 + 12 = 29,2; All sides / All profits
%   16 + 1,4 + 12 = 29,4.
%   „New order“: Zeile (Plus + Abstand + Text) mittig im Inhalt x20,8..99,2 (Mitte 60);
%   sie beginnt beim Plus (x22,3), endet also bei 60 + (60 - 22,3) = 97,7 = Textende.
%
% Vollständige Fassung (Abweichungen vom Original, nur Weglassen):
%   Bildspalte (x10..80, 70 px) entfällt: alle Spalten ab ID um -70 px verschoben, die
%   Zeilen beginnen wieder mit Radius 3 px. Zahnrad- und Schließen-Spalte (x1544..1694,
%   150 px) entfallen; Zeilenende wieder Radius 3 px. Tabelle x10..1474, Block 1484 px
%   breit (1704 - 70 - 150). Die Reiter bleiben bündig mit dem rechten Tabellenrand, die
%   Leiste bleibt an ihrer Stelle oben links (x10..575).
%   unten  Schnitt 10 px unter der zweiten Zeile (y145; leere Fläche bis y271 entfällt),
%          Ecken wieder Radius 7 px.
%   Weggelassen: Griff zum Einklappen (a.hide, Mitte oben, Icon U+E90E fehlt in
%   icons.tikz), Instrument-Bild in der Bildspalte und Pfeil hinter dem Instrument (PNG,
%   kein Vektor; die Fläche bleibt leer: Text um 5,5 px nach links wie im Original),
%   Sortierpfeile (padding-right 15 px der Köpfe; Inhalt nicht gemessen), Pfeil der
%   Auswahlfelder All sides / All profits (ng-arrow, nicht gemessen).
%   Plus (U+E91A) und Lupe (U+E915) fehlen in icons.tikz: Fläche leer; sie werden
%   automatisch gezeichnet, sobald \iconplus bzw. \iconsearch existieren.
%   NICHT GEMESSEN, vorläufig: Farbe des Platzhalters „Search“ (::placeholder) – hier
%   rgb(168,187,190) wie Köpfe und Lupe (\definecolor{suchtext} unten).
%
% Das Dokument ist zahlenfrei: Werte sind Wörter, ohne „€“ (wie bar.tex). ID „Nummer“,
% Instrument „Instrument“, Type „Buy“, Amount „Menge“, Trade Volume „Wert“, Open Rate
% „Kurs“, Open Time „Zeitpunkt“, S/L, T/P und Commission „—“, Swap „Betrag“,
% Total Profit „Ergebnis“ (grün/rot), Margin „Betrag“. Reiter ohne Zähler; ihre Breite
% rechnet wie der Browser (16 + Text + 16 px). Probe der Textbreiten mit den gemessenen
% Werten: siehe \typeout POSITIONEN-VOLL im Log.
%
% Bei width=257mm: 1 px = 0,1732 mm, Bild 257 x 25,1 mm, Schrift 12 px = 5,9 pt.
%
% Marken (Schalter \ifmarken, Standard: an; ohne Marken bauen mit
%   xelatex '\def\ohnemarken{}% Tabelle der offenen Positionen am Desktop (Reiter Active, blaues Design) – Vektor-Nachbau
% VOLLSTÄNDIGE Fassung ohne Bruchkante, für eine Querformat-Seite (Satzbreite ca. 257 mm).
% Gegenstück zum Ausschnitt positionen-desktop.tex (gleiche Technik, Farben, Platzhalter).
% Grundlage: DOM-Messung originals/messungen/positionen-desktop.json (CSS-px, Fenster
% 2040 x 986; anonymisiert: Ziffern = 0, Instrumente = XXXXXX). Koordinaten unten in
% CSS-px, y nach unten, Ursprung = linke obere Ecke des Blocks.
%
% Gemessen (Original-x):
%   .maintable   1704 x 271, Fläche rgb(4,39,60), Radius 7 px, Inter 400, 12 px / 14 px
%   Leiste oben links (Lage im Block, unabhängig von den Spalten):
%     a.button.bordered  „New order“ x10 y10,2 100 x 29,6, Rand 0,8 px rgb(83,188,81),
%                        Radius 5 px, Fläche frei, padding 6px 10px, Inter 600 weiß,
%                        Zeile 16 px, text-align center
%       i.icon-plus      x22,3 y17,8 12 x 12, icomoon 12 px, U+E91A, weiß
%     input.form-control x120 y10 145 x 30, Fläche und Rand rgb(0,19,40), Radius 4 px,
%                        padding 6px 30px 6px 10px, Zeile 16 px; Platzhalter „Search“
%       i.icon-search    x238 y17,5 15 x 15, icomoon 15 px, U+E915, rgb(168,187,190)
%     ng-select-container x275 / x430 y10 145 x 30, Fläche und Rand rgb(0,19,40),
%                        Radius 4 px; ng-value x275,8 / x430,8, padding 6px 10px,
%                        Zeile 18 px, weiß: „All sides“ (67,4), „All profits“ (73,8)
%   ul.nav-tabs  x1436,3 y10 257,7 x 30, Fläche rgb(0,19,40), Radius 3 px; endet bündig
%                mit der Tabelle (x1694 = 1704 - 10)
%     a.nav-link  flex, padding 0 16px, Text mittig, weiß, capitalize:
%                „Active (0)“ 87 (aktiv: Fläche rgb(14,80,114), Radius 3 0 0 3),
%                „Pending (0)“ 98,1, „History“ 72,6 (Radius 0 3 3 0)
%   th           y47, 28,8 hoch, padding 4px 10px, Text rgb(168,187,190), capitalize;
%                span (inline-block, 20 px Zeile) mit padding-right 15 px; Text beginnt am
%                gemessenen span-x: ID 102,5, Instrument 177, Type 298, Amount 383,8,
%                Trade volume 493,5, Open Rate 640, Open Time 787,5, S/L 941, T/P 1030,9,
%                Swap 1113,6, Commission 1207,4, Total Profit 1344,2, Margin 1467,9
%   td           Zeilen y75,8 / y105,4, je 29,6 hoch, border-top 1,6 px rgb(4,39,60),
%                Fläche rgba(14,80,114,.1), Text mittig (padding 4px 10px); erste Zelle
%                Radius 3 0 0 3, letzte 0 3 3 0
%                Spalten: (Bild 10 w70) | ID 80 w71,8 | Instrument 151,8 w126,2 |
%                Type 278 w83,2 | Amount 361,2 w104,3 | Trade volume 465,5 w149,5 |
%                Open Rate 615 w124,2 | Open Time 739,2 w173,8 | S/L 912,9 w90 |
%                T/P 1002,9 w90 | Swap 1093 w87,8 | Commission 1180,8 w138,1 |
%                Total Profit 1318,9 w127,2 | Margin 1446 w98 | (Zahnrad 1544 w50) |
%                (Schließen 1594 w100)
%   Total Profit grün rgb(83,188,81) (color-green-main) bzw. rot rgb(249,64,86)
%                (color-red-main)
% Grundlinien (Chrome, Gerätepixel 0,8 px: Ascent 12, Descent 3,2 px; icomoon 12 px:
% Ascent 11,2, Descent 0,8 px): Zeilenoberkante + (Zeilenhöhe - 15,2) / 2 + 12.
%   Kopf 65,4; Zeilen 95,0 und 124,6 (belegt durch €-span und Bildmitte, siehe
%   positionen-desktop.tex); Reiter 29,2 (mittig angenommen, nicht gemessen);
%   New order 29,0 (belegt: Oberkante icon-plus 17,8 + Ascent 11,2; der halbe Durchschuss
%   0,4 px fällt weg); Search 10 + 0,8 + 6 + 0,4 + 12 = 29,2; All sides / All profits
%   16 + 1,4 + 12 = 29,4.
%   „New order“: Zeile (Plus + Abstand + Text) mittig im Inhalt x20,8..99,2 (Mitte 60);
%   sie beginnt beim Plus (x22,3), endet also bei 60 + (60 - 22,3) = 97,7 = Textende.
%
% Vollständige Fassung (Abweichungen vom Original, nur Weglassen):
%   Bildspalte (x10..80, 70 px) entfällt: alle Spalten ab ID um -70 px verschoben, die
%   Zeilen beginnen wieder mit Radius 3 px. Zahnrad- und Schließen-Spalte (x1544..1694,
%   150 px) entfallen; Zeilenende wieder Radius 3 px. Tabelle x10..1474, Block 1484 px
%   breit (1704 - 70 - 150). Die Reiter bleiben bündig mit dem rechten Tabellenrand, die
%   Leiste bleibt an ihrer Stelle oben links (x10..575).
%   unten  Schnitt 10 px unter der zweiten Zeile (y145; leere Fläche bis y271 entfällt),
%          Ecken wieder Radius 7 px.
%   Weggelassen: Griff zum Einklappen (a.hide, Mitte oben, Icon U+E90E fehlt in
%   icons.tikz), Instrument-Bild in der Bildspalte und Pfeil hinter dem Instrument (PNG,
%   kein Vektor; die Fläche bleibt leer: Text um 5,5 px nach links wie im Original),
%   Sortierpfeile (padding-right 15 px der Köpfe; Inhalt nicht gemessen), Pfeil der
%   Auswahlfelder All sides / All profits (ng-arrow, nicht gemessen).
%   Plus (U+E91A) und Lupe (U+E915) fehlen in icons.tikz: Fläche leer; sie werden
%   automatisch gezeichnet, sobald \iconplus bzw. \iconsearch existieren.
%   NICHT GEMESSEN, vorläufig: Farbe des Platzhalters „Search“ (::placeholder) – hier
%   rgb(168,187,190) wie Köpfe und Lupe (\definecolor{suchtext} unten).
%
% Das Dokument ist zahlenfrei: Werte sind Wörter, ohne „€“ (wie bar.tex). ID „Nummer“,
% Instrument „Instrument“, Type „Buy“, Amount „Menge“, Trade Volume „Wert“, Open Rate
% „Kurs“, Open Time „Zeitpunkt“, S/L, T/P und Commission „—“, Swap „Betrag“,
% Total Profit „Ergebnis“ (grün/rot), Margin „Betrag“. Reiter ohne Zähler; ihre Breite
% rechnet wie der Browser (16 + Text + 16 px). Probe der Textbreiten mit den gemessenen
% Werten: siehe \typeout POSITIONEN-VOLL im Log.
%
% Bei width=257mm: 1 px = 0,1732 mm, Bild 257 x 25,1 mm, Schrift 12 px = 5,9 pt.
%
% Marken (Schalter \ifmarken, Standard: an; ohne Marken bauen mit
%   xelatex '\def\ohnemarken{}% Tabelle der offenen Positionen am Desktop (Reiter Active, blaues Design) – Vektor-Nachbau
% VOLLSTÄNDIGE Fassung ohne Bruchkante, für eine Querformat-Seite (Satzbreite ca. 257 mm).
% Gegenstück zum Ausschnitt positionen-desktop.tex (gleiche Technik, Farben, Platzhalter).
% Grundlage: DOM-Messung originals/messungen/positionen-desktop.json (CSS-px, Fenster
% 2040 x 986; anonymisiert: Ziffern = 0, Instrumente = XXXXXX). Koordinaten unten in
% CSS-px, y nach unten, Ursprung = linke obere Ecke des Blocks.
%
% Gemessen (Original-x):
%   .maintable   1704 x 271, Fläche rgb(4,39,60), Radius 7 px, Inter 400, 12 px / 14 px
%   Leiste oben links (Lage im Block, unabhängig von den Spalten):
%     a.button.bordered  „New order“ x10 y10,2 100 x 29,6, Rand 0,8 px rgb(83,188,81),
%                        Radius 5 px, Fläche frei, padding 6px 10px, Inter 600 weiß,
%                        Zeile 16 px, text-align center
%       i.icon-plus      x22,3 y17,8 12 x 12, icomoon 12 px, U+E91A, weiß
%     input.form-control x120 y10 145 x 30, Fläche und Rand rgb(0,19,40), Radius 4 px,
%                        padding 6px 30px 6px 10px, Zeile 16 px; Platzhalter „Search“
%       i.icon-search    x238 y17,5 15 x 15, icomoon 15 px, U+E915, rgb(168,187,190)
%     ng-select-container x275 / x430 y10 145 x 30, Fläche und Rand rgb(0,19,40),
%                        Radius 4 px; ng-value x275,8 / x430,8, padding 6px 10px,
%                        Zeile 18 px, weiß: „All sides“ (67,4), „All profits“ (73,8)
%   ul.nav-tabs  x1436,3 y10 257,7 x 30, Fläche rgb(0,19,40), Radius 3 px; endet bündig
%                mit der Tabelle (x1694 = 1704 - 10)
%     a.nav-link  flex, padding 0 16px, Text mittig, weiß, capitalize:
%                „Active (0)“ 87 (aktiv: Fläche rgb(14,80,114), Radius 3 0 0 3),
%                „Pending (0)“ 98,1, „History“ 72,6 (Radius 0 3 3 0)
%   th           y47, 28,8 hoch, padding 4px 10px, Text rgb(168,187,190), capitalize;
%                span (inline-block, 20 px Zeile) mit padding-right 15 px; Text beginnt am
%                gemessenen span-x: ID 102,5, Instrument 177, Type 298, Amount 383,8,
%                Trade volume 493,5, Open Rate 640, Open Time 787,5, S/L 941, T/P 1030,9,
%                Swap 1113,6, Commission 1207,4, Total Profit 1344,2, Margin 1467,9
%   td           Zeilen y75,8 / y105,4, je 29,6 hoch, border-top 1,6 px rgb(4,39,60),
%                Fläche rgba(14,80,114,.1), Text mittig (padding 4px 10px); erste Zelle
%                Radius 3 0 0 3, letzte 0 3 3 0
%                Spalten: (Bild 10 w70) | ID 80 w71,8 | Instrument 151,8 w126,2 |
%                Type 278 w83,2 | Amount 361,2 w104,3 | Trade volume 465,5 w149,5 |
%                Open Rate 615 w124,2 | Open Time 739,2 w173,8 | S/L 912,9 w90 |
%                T/P 1002,9 w90 | Swap 1093 w87,8 | Commission 1180,8 w138,1 |
%                Total Profit 1318,9 w127,2 | Margin 1446 w98 | (Zahnrad 1544 w50) |
%                (Schließen 1594 w100)
%   Total Profit grün rgb(83,188,81) (color-green-main) bzw. rot rgb(249,64,86)
%                (color-red-main)
% Grundlinien (Chrome, Gerätepixel 0,8 px: Ascent 12, Descent 3,2 px; icomoon 12 px:
% Ascent 11,2, Descent 0,8 px): Zeilenoberkante + (Zeilenhöhe - 15,2) / 2 + 12.
%   Kopf 65,4; Zeilen 95,0 und 124,6 (belegt durch €-span und Bildmitte, siehe
%   positionen-desktop.tex); Reiter 29,2 (mittig angenommen, nicht gemessen);
%   New order 29,0 (belegt: Oberkante icon-plus 17,8 + Ascent 11,2; der halbe Durchschuss
%   0,4 px fällt weg); Search 10 + 0,8 + 6 + 0,4 + 12 = 29,2; All sides / All profits
%   16 + 1,4 + 12 = 29,4.
%   „New order“: Zeile (Plus + Abstand + Text) mittig im Inhalt x20,8..99,2 (Mitte 60);
%   sie beginnt beim Plus (x22,3), endet also bei 60 + (60 - 22,3) = 97,7 = Textende.
%
% Vollständige Fassung (Abweichungen vom Original, nur Weglassen):
%   Bildspalte (x10..80, 70 px) entfällt: alle Spalten ab ID um -70 px verschoben, die
%   Zeilen beginnen wieder mit Radius 3 px. Zahnrad- und Schließen-Spalte (x1544..1694,
%   150 px) entfallen; Zeilenende wieder Radius 3 px. Tabelle x10..1474, Block 1484 px
%   breit (1704 - 70 - 150). Die Reiter bleiben bündig mit dem rechten Tabellenrand, die
%   Leiste bleibt an ihrer Stelle oben links (x10..575).
%   unten  Schnitt 10 px unter der zweiten Zeile (y145; leere Fläche bis y271 entfällt),
%          Ecken wieder Radius 7 px.
%   Weggelassen: Griff zum Einklappen (a.hide, Mitte oben, Icon U+E90E fehlt in
%   icons.tikz), Instrument-Bild in der Bildspalte und Pfeil hinter dem Instrument (PNG,
%   kein Vektor; die Fläche bleibt leer: Text um 5,5 px nach links wie im Original),
%   Sortierpfeile (padding-right 15 px der Köpfe; Inhalt nicht gemessen), Pfeil der
%   Auswahlfelder All sides / All profits (ng-arrow, nicht gemessen).
%   Plus (U+E91A) und Lupe (U+E915) fehlen in icons.tikz: Fläche leer; sie werden
%   automatisch gezeichnet, sobald \iconplus bzw. \iconsearch existieren.
%   NICHT GEMESSEN, vorläufig: Farbe des Platzhalters „Search“ (::placeholder) – hier
%   rgb(168,187,190) wie Köpfe und Lupe (\definecolor{suchtext} unten).
%
% Das Dokument ist zahlenfrei: Werte sind Wörter, ohne „€“ (wie bar.tex). ID „Nummer“,
% Instrument „Instrument“, Type „Buy“, Amount „Menge“, Trade Volume „Wert“, Open Rate
% „Kurs“, Open Time „Zeitpunkt“, S/L, T/P und Commission „—“, Swap „Betrag“,
% Total Profit „Ergebnis“ (grün/rot), Margin „Betrag“. Reiter ohne Zähler; ihre Breite
% rechnet wie der Browser (16 + Text + 16 px). Probe der Textbreiten mit den gemessenen
% Werten: siehe \typeout POSITIONEN-VOLL im Log.
%
% Bei width=257mm: 1 px = 0,1732 mm, Bild 257 x 25,1 mm, Schrift 12 px = 5,9 pt.
%
% Marken (Schalter \ifmarken, Standard: an; ohne Marken bauen mit
%   xelatex '\def\ohnemarken{}% Tabelle der offenen Positionen am Desktop (Reiter Active, blaues Design) – Vektor-Nachbau
% VOLLSTÄNDIGE Fassung ohne Bruchkante, für eine Querformat-Seite (Satzbreite ca. 257 mm).
% Gegenstück zum Ausschnitt positionen-desktop.tex (gleiche Technik, Farben, Platzhalter).
% Grundlage: DOM-Messung originals/messungen/positionen-desktop.json (CSS-px, Fenster
% 2040 x 986; anonymisiert: Ziffern = 0, Instrumente = XXXXXX). Koordinaten unten in
% CSS-px, y nach unten, Ursprung = linke obere Ecke des Blocks.
%
% Gemessen (Original-x):
%   .maintable   1704 x 271, Fläche rgb(4,39,60), Radius 7 px, Inter 400, 12 px / 14 px
%   Leiste oben links (Lage im Block, unabhängig von den Spalten):
%     a.button.bordered  „New order“ x10 y10,2 100 x 29,6, Rand 0,8 px rgb(83,188,81),
%                        Radius 5 px, Fläche frei, padding 6px 10px, Inter 600 weiß,
%                        Zeile 16 px, text-align center
%       i.icon-plus      x22,3 y17,8 12 x 12, icomoon 12 px, U+E91A, weiß
%     input.form-control x120 y10 145 x 30, Fläche und Rand rgb(0,19,40), Radius 4 px,
%                        padding 6px 30px 6px 10px, Zeile 16 px; Platzhalter „Search“
%       i.icon-search    x238 y17,5 15 x 15, icomoon 15 px, U+E915, rgb(168,187,190)
%     ng-select-container x275 / x430 y10 145 x 30, Fläche und Rand rgb(0,19,40),
%                        Radius 4 px; ng-value x275,8 / x430,8, padding 6px 10px,
%                        Zeile 18 px, weiß: „All sides“ (67,4), „All profits“ (73,8)
%   ul.nav-tabs  x1436,3 y10 257,7 x 30, Fläche rgb(0,19,40), Radius 3 px; endet bündig
%                mit der Tabelle (x1694 = 1704 - 10)
%     a.nav-link  flex, padding 0 16px, Text mittig, weiß, capitalize:
%                „Active (0)“ 87 (aktiv: Fläche rgb(14,80,114), Radius 3 0 0 3),
%                „Pending (0)“ 98,1, „History“ 72,6 (Radius 0 3 3 0)
%   th           y47, 28,8 hoch, padding 4px 10px, Text rgb(168,187,190), capitalize;
%                span (inline-block, 20 px Zeile) mit padding-right 15 px; Text beginnt am
%                gemessenen span-x: ID 102,5, Instrument 177, Type 298, Amount 383,8,
%                Trade volume 493,5, Open Rate 640, Open Time 787,5, S/L 941, T/P 1030,9,
%                Swap 1113,6, Commission 1207,4, Total Profit 1344,2, Margin 1467,9
%   td           Zeilen y75,8 / y105,4, je 29,6 hoch, border-top 1,6 px rgb(4,39,60),
%                Fläche rgba(14,80,114,.1), Text mittig (padding 4px 10px); erste Zelle
%                Radius 3 0 0 3, letzte 0 3 3 0
%                Spalten: (Bild 10 w70) | ID 80 w71,8 | Instrument 151,8 w126,2 |
%                Type 278 w83,2 | Amount 361,2 w104,3 | Trade volume 465,5 w149,5 |
%                Open Rate 615 w124,2 | Open Time 739,2 w173,8 | S/L 912,9 w90 |
%                T/P 1002,9 w90 | Swap 1093 w87,8 | Commission 1180,8 w138,1 |
%                Total Profit 1318,9 w127,2 | Margin 1446 w98 | (Zahnrad 1544 w50) |
%                (Schließen 1594 w100)
%   Total Profit grün rgb(83,188,81) (color-green-main) bzw. rot rgb(249,64,86)
%                (color-red-main)
% Grundlinien (Chrome, Gerätepixel 0,8 px: Ascent 12, Descent 3,2 px; icomoon 12 px:
% Ascent 11,2, Descent 0,8 px): Zeilenoberkante + (Zeilenhöhe - 15,2) / 2 + 12.
%   Kopf 65,4; Zeilen 95,0 und 124,6 (belegt durch €-span und Bildmitte, siehe
%   positionen-desktop.tex); Reiter 29,2 (mittig angenommen, nicht gemessen);
%   New order 29,0 (belegt: Oberkante icon-plus 17,8 + Ascent 11,2; der halbe Durchschuss
%   0,4 px fällt weg); Search 10 + 0,8 + 6 + 0,4 + 12 = 29,2; All sides / All profits
%   16 + 1,4 + 12 = 29,4.
%   „New order“: Zeile (Plus + Abstand + Text) mittig im Inhalt x20,8..99,2 (Mitte 60);
%   sie beginnt beim Plus (x22,3), endet also bei 60 + (60 - 22,3) = 97,7 = Textende.
%
% Vollständige Fassung (Abweichungen vom Original, nur Weglassen):
%   Bildspalte (x10..80, 70 px) entfällt: alle Spalten ab ID um -70 px verschoben, die
%   Zeilen beginnen wieder mit Radius 3 px. Zahnrad- und Schließen-Spalte (x1544..1694,
%   150 px) entfallen; Zeilenende wieder Radius 3 px. Tabelle x10..1474, Block 1484 px
%   breit (1704 - 70 - 150). Die Reiter bleiben bündig mit dem rechten Tabellenrand, die
%   Leiste bleibt an ihrer Stelle oben links (x10..575).
%   unten  Schnitt 10 px unter der zweiten Zeile (y145; leere Fläche bis y271 entfällt),
%          Ecken wieder Radius 7 px.
%   Weggelassen: Griff zum Einklappen (a.hide, Mitte oben, Icon U+E90E fehlt in
%   icons.tikz), Instrument-Bild in der Bildspalte und Pfeil hinter dem Instrument (PNG,
%   kein Vektor; die Fläche bleibt leer: Text um 5,5 px nach links wie im Original),
%   Sortierpfeile (padding-right 15 px der Köpfe; Inhalt nicht gemessen), Pfeil der
%   Auswahlfelder All sides / All profits (ng-arrow, nicht gemessen).
%   Plus (U+E91A) und Lupe (U+E915) fehlen in icons.tikz: Fläche leer; sie werden
%   automatisch gezeichnet, sobald \iconplus bzw. \iconsearch existieren.
%   NICHT GEMESSEN, vorläufig: Farbe des Platzhalters „Search“ (::placeholder) – hier
%   rgb(168,187,190) wie Köpfe und Lupe (\definecolor{suchtext} unten).
%
% Das Dokument ist zahlenfrei: Werte sind Wörter, ohne „€“ (wie bar.tex). ID „Nummer“,
% Instrument „Instrument“, Type „Buy“, Amount „Menge“, Trade Volume „Wert“, Open Rate
% „Kurs“, Open Time „Zeitpunkt“, S/L, T/P und Commission „—“, Swap „Betrag“,
% Total Profit „Ergebnis“ (grün/rot), Margin „Betrag“. Reiter ohne Zähler; ihre Breite
% rechnet wie der Browser (16 + Text + 16 px). Probe der Textbreiten mit den gemessenen
% Werten: siehe \typeout POSITIONEN-VOLL im Log.
%
% Bei width=257mm: 1 px = 0,1732 mm, Bild 257 x 25,1 mm, Schrift 12 px = 5,9 pt.
%
% Marken (Schalter \ifmarken, Standard: an; ohne Marken bauen mit
%   xelatex '\def\ohnemarken{}\input{positionen-desktop-voll.tex}'):
% Stil wie info-desktop.tex (weißer Kreis 17 px, Ziffer Inter 600 in #04273C), mittig auf
% der Kopfzeile, 3 px vor dem Spaltennamen: 1 Swap, 2 Total Profit, 3 Margin.
\documentclass[border=0pt]{standalone}
\usepackage{fontspec}
\usepackage{tikz}
\setmainfont{Inter}% Inter 4.0 (Paket fonts-inter), wie bar.tex
\newfontfamily\halbfett{Inter SemiBold}% New order, Ziffern der Marken (wie panel.pdf)

\newlength\px \setlength\px{0.75bp}% 1 CSS-px
\newcommand\schrift{\fontsize{9bp}{10.5bp}\selectfont}% 12 px / line-height 14 px
\definecolor{block}{RGB}{4,39,60}%      .maintable
\definecolor{reiter}{RGB}{0,19,40}%     ul.nav-tabs, input, ng-select-container
\definecolor{aktiv}{RGB}{14,80,114}%    a.nav-link.active; Zeilen mit 10 % Deckung
\definecolor{label}{RGB}{168,187,190}%  th, i.icon-search
\definecolor{gewinn}{RGB}{83,188,81}%   color-green-main, Rand von New order
\definecolor{verlust}{RGB}{249,64,86}%  color-red-main
\definecolor{suchtext}{RGB}{168,187,190}% Platzhalter „Search“ – NICHT GEMESSEN (vorläufig)
\definecolor{markenziffer}{HTML}{04273C}% Ziffer der Marken (wie panel.pdf)

\input{icons.tikz}

\newif\ifmarken \markentrue
\ifdefined\ohnemarken \markenfalse\fi

\newcommand\breite{1484}% 1704 - 70 (Bildspalte) - 150 (Zahnrad, Schließen)
\newcommand\hoehe{145}%   zweite Zeile endet bei 135, + 10 px
\newcommand\links{-70}%   Verschiebung der Tabellenspalten (ID beginnt bei x10)
\newcommand\rand{1474}%   rechter Tabellenrand (\breite - 10)

% Textbreiten in px (vor dem Bild messen; in TikZ ist \nullfont aktiv)
\newlength\tw
\makeatletter
\newcommand\breitepx[2]{\settowidth\tw{\schrift #2}%
  \edef#1{\strip@pt\dimexpr\tw*4/3*7200/7227\relax}}
\newcommand\probe[2]{\breitepx\probewert{#1}%
  \typeout{POSITIONEN-VOLL: '#1' \probewert\space px (gemessen #2)}}
\makeatother
\probe{Active (0)}{55,0}\probe{Pending (0)}{66,1}\probe{History}{40,6}
\probe{ID}{11,9}\probe{Instrument}{60,7}\probe{Type}{28,1}\probe{Amount}{44,1}
\probe{Trade Volume}{78,4}\probe{Open Rate}{59,2}\probe{Open Time}{62,2}
\probe{S/L}{18,8}\probe{T/P}{19,1}\probe{Swap}{31,6}
\probe{Commission}{69,8}\probe{Total Profit}{61,5}\probe{Margin}{39,3}
\probe{All sides}{47,4}\probe{All profits}{53,8}
\breitepx\probewert{Search}\typeout{POSITIONEN-VOLL: 'Search' \probewert\space px (nicht gemessen)}
\breitepx\wn{\halbfett New order}\pgfmathsetmacro\probewert{97.7-\wn-22.3-12}
\typeout{POSITIONEN-VOLL: 'New order' (600) \wn\space px; Abstand Plus-Text \probewert\space px}
\breitepx\wa{Active}\breitepx\wp{Pending}\breitepx\wh{History}
\pgfmathsetmacro\xh{\rand-(\wh+32)}% Reiter rechtsbündig am Tabellenrand
\pgfmathsetmacro\xp{\xh-(\wp+32)}
\pgfmathsetmacro\xa{\xp-(\wa+32)}
\typeout{POSITIONEN-VOLL: Reiter ab x \xa, \xp, \xh\space bis \rand\space px}

% \flaeche{<Zeilenoberkante>}: getönte Zeile, Zellen-Fläche unter dem 1,6-px-Rand;
% außen Radius 3 px (oben innen 3 x 1,4 px)
\newcommand\flaeche[1]{%
  \fill[aktiv,fill opacity=0.1] (\breite/2,#1+1.6) -- (13,#1+1.6)
    arc[start angle=270,end angle=180,x radius=3,y radius=1.4] -- (10,#1+26.6)
    arc[start angle=180,end angle=90,radius=3] -- (\rand-3,#1+29.6)
    arc[start angle=90,end angle=0,radius=3] -- (\rand,#1+3)
    arc[start angle=0,end angle=-90,x radius=3,y radius=1.4] -- cycle;}

% \kopf{<Original-x des span>}{<Text>} / \zelle{<Mitte>}{<y>}{<Farbe>}{<Text>}
\newcommand\kopf[2]{\node[anchor=base west,text=label] at (#1,65.4) {#2};}
\newcommand\zelle[4]{\node[anchor=base,text=#3] at (#1,#2) {#4};}

% \icon{<x>}{<y Oberkante>}{<Größe px>}{<Farbe>}{<Icon>}: Box Größe x Größe
% (icons.tikz: 24 px, Ursprung unten links, y nach oben)
\newcommand\icon[5]{%
  \begin{scope}[shift={(#1,#2+#3)},x=#3/32*1bp,y=#3/32*1bp,color=#4]#5\end{scope}}

% \marke{<Original-x des Spaltennamens>}{<Nummer>}: Kreis 17 px, Mitte 3 px + Radius davor
\newcommand\marke[2]{%
  \fill[white] (#1-11.5,61) circle[radius=8.5];
  \node[anchor=base,text=markenziffer,font=\fontsize{6.954bp}{8bp}\halbfett]
    at (#1-11.5,61+3.091) {#2};}% 12 px * 17/22

\begin{document}%
\begin{tikzpicture}[x=\px,y={(0,-\px)},font=\schrift,
    every node/.style={inner sep=0,outer sep=0}]
  \path[use as bounding box] (0,0) rectangle (\breite,\hoehe);
  % Block und Zeilen
  \fill[block,rounded corners=7\px] (0,0) rectangle (\breite,\hoehe);
  \flaeche{75.8}\flaeche{105.4}
  % Leiste oben links (Lage wie gemessen)
  \draw[gewinn,line width=0.8\px,rounded corners=4.6\px]
    (10.4,10.6) rectangle (109.6,39.4);% Rand 0,8 px innen, Radius 5 px außen
  \ifdefined\iconplus \icon{22.3}{17.8}{12}{white}{\iconplus}\fi
  \node[anchor=base east,text=white,font=\schrift\halbfett] at (97.7,29.0) {New order};
  \fill[reiter,rounded corners=4\px] (120,10) rectangle (265,40);
  \node[anchor=base west,text=suchtext] at (130.8,29.2) {Search};
  \ifdefined\iconsearch \icon{238}{17.5}{15}{label}{\iconsearch}\fi
  \fill[reiter,rounded corners=4\px] (275,10) rectangle (420,40);
  \node[anchor=base west,text=white] at (285.8,29.4) {All sides};
  \fill[reiter,rounded corners=4\px] (430,10) rectangle (575,40);
  \node[anchor=base west,text=white] at (440.8,29.4) {All profits};
  % Reiter (ohne Zähler), bündig mit dem rechten Tabellenrand
  \fill[reiter,rounded corners=3\px] (\xa,10) rectangle (\rand,40);
  \fill[aktiv] (\xp,10) {[rounded corners=3\px] -- (\xa,10) -- (\xa,40)} -- (\xp,40) -- cycle;
  \node[anchor=base,text=white] at ({\xa+(\wa+32)/2},29.2) {Active};
  \node[anchor=base,text=white] at ({\xp+(\wp+32)/2},29.2) {Pending};
  \node[anchor=base,text=white] at ({\xh+(\wh+32)/2},29.2) {History};
  % Tabelle (Original-x, um die Bildspalte nach links verschoben)
  \begin{scope}[xshift=\links\px]
    \kopf{102.5}{ID}\kopf{177}{Instrument}\kopf{298}{Type}\kopf{383.8}{Amount}
    \kopf{493.5}{Trade Volume}\kopf{640}{Open Rate}\kopf{787.5}{Open Time}
    \kopf{941}{S/L}\kopf{1030.9}{T/P}\kopf{1113.6}{Swap}\kopf{1207.4}{Commission}
    \kopf{1344.2}{Total Profit}\kopf{1467.9}{Margin}
    \foreach \y/\farbe in {95.0/gewinn,124.6/verlust} {
      \zelle{115.9}{\y}{white}{Nummer}
      \zelle{209.4}{\y}{white}{Instrument}% Mitte 214,9 - 5,5 (Pfeil-Bild weggelassen)
      \zelle{319.6}{\y}{white}{Buy}
      \zelle{413.35}{\y}{white}{Menge}
      \zelle{540.25}{\y}{white}{Wert}
      \zelle{677.1}{\y}{white}{Kurs}
      \zelle{826.1}{\y}{white}{Zeitpunkt}
      \zelle{957.9}{\y}{white}{—}
      \zelle{1047.9}{\y}{white}{—}
      \zelle{1136.9}{\y}{white}{Betrag}
      \zelle{1249.85}{\y}{white}{—}
      \zelle{1382.5}{\y}{\farbe}{Ergebnis}
      \zelle{1495}{\y}{white}{Betrag}}
    \ifmarken
      \marke{1113.6}{1}
      \marke{1344.2}{2}
      \marke{1467.9}{3}
    \fi
  \end{scope}
\end{tikzpicture}%
\end{document}
'):
% Stil wie info-desktop.tex (weißer Kreis 17 px, Ziffer Inter 600 in #04273C), mittig auf
% der Kopfzeile, 3 px vor dem Spaltennamen: 1 Swap, 2 Total Profit, 3 Margin.
\documentclass[border=0pt]{standalone}
\usepackage{fontspec}
\usepackage{tikz}
\setmainfont{Inter}% Inter 4.0 (Paket fonts-inter), wie bar.tex
\newfontfamily\halbfett{Inter SemiBold}% New order, Ziffern der Marken (wie panel.pdf)

\newlength\px \setlength\px{0.75bp}% 1 CSS-px
\newcommand\schrift{\fontsize{9bp}{10.5bp}\selectfont}% 12 px / line-height 14 px
\definecolor{block}{RGB}{4,39,60}%      .maintable
\definecolor{reiter}{RGB}{0,19,40}%     ul.nav-tabs, input, ng-select-container
\definecolor{aktiv}{RGB}{14,80,114}%    a.nav-link.active; Zeilen mit 10 % Deckung
\definecolor{label}{RGB}{168,187,190}%  th, i.icon-search
\definecolor{gewinn}{RGB}{83,188,81}%   color-green-main, Rand von New order
\definecolor{verlust}{RGB}{249,64,86}%  color-red-main
\definecolor{suchtext}{RGB}{168,187,190}% Platzhalter „Search“ – NICHT GEMESSEN (vorläufig)
\definecolor{markenziffer}{HTML}{04273C}% Ziffer der Marken (wie panel.pdf)

% Icons als TikZ-Pfade – erzeugt von tools/icons_aus_font.py, nicht von Hand ändern.
% Schrift: 230d1684-icomoon-1-icomoon.b9abfeda36b81d3e.ttf (icomoon, Version 1.0, TrueType), sha256 951e3a5eca89112f…
% Herkunft: Webfont „icomoon“ der Plattform (icomoon.b9abfeda36b81d3e.ttf; icomoon.47a83e5a26e26197.woff hat dieselben Tabellen und Glyphen), vom Nutzer hochgeladen am 02.10.2026. Die Schriftdateien liegen NICHT im Repo (Lizenz unklar), nur diese Pfade.
% unitsPerEm 1024; Metriken hhea: Ascent 960, Descent -64.
% Box wie das <i>-Element im Browser: font-size 24px, line-height 24px
% -> Box 24 x 24 px, Ursprung unten links, y nach oben;
%    Grundlinie 1.500 px über der Boxunterkante. 1 Einheit = 1 CSS-px.
% Füllregel nonzero (wie im Browser). Farbe kommt von außen (color=…).
% Gebrauch: \begin{scope}[shift={(<x>,<y>)},x=\px,y=\px,color=<Farbe>]\iconbell\end{scope}

% U+E938 uniE938 – icon-notifications (Glocke); Vorschub 24.000 px; Umriss x 1.99–22.01, y 1.99–22.01 px; 3 Kontur(en)
\newcommand\iconbell{\fill[nonzero rule]
    (19.992,7.711) -- (22.008,7.711) -- (22.008,5.812) -- (1.992,5.812) -- (1.992,7.711) -- (4.008,7.711) -- (4.008,14.391) .. controls (4.008,15.391) and (4.211,16.359) .. (4.617,17.297) .. controls (5.023,18.234) and (5.602,19.055) .. (6.352,19.758) .. controls (7.102,20.477) and (7.965,21.031) .. (8.941,21.422) .. controls (9.918,21.812) and (10.938,22.008) .. (12,22.008) .. controls (13.062,22.008) and (14.082,21.812) .. (15.059,21.422) .. controls (16.035,21.031) and (16.898,20.477) .. (17.648,19.758) .. controls (18.398,19.055) and (18.977,18.234) .. (19.383,17.297) .. controls (19.789,16.359) and (19.992,15.391) .. (19.992,14.391) -- (19.992,7.711) -- cycle
    (18,7.711) -- (18,14.391) .. controls (18,15.141) and (17.848,15.867) .. (17.543,16.57) .. controls (17.238,17.273) and (16.805,17.891) .. (16.242,18.422) .. controls (15.68,18.953) and (15.031,19.363) .. (14.297,19.652) .. controls (13.562,19.941) and (12.797,20.086) .. (12,20.086) .. controls (11.203,20.086) and (10.438,19.941) .. (9.703,19.652) .. controls (8.969,19.363) and (8.32,18.953) .. (7.758,18.422) .. controls (7.195,17.891) and (6.762,17.273) .. (6.457,16.57) .. controls (6.152,15.867) and (6,15.141) .. (6,14.391) -- (6,7.711) -- cycle
    (9,3.891) -- (15,3.891) -- (15,1.992) -- (9,1.992) -- cycle;}

% U+E937 uniE937 – icon-settings (Zahnrad); Vorschub 24.000 px; Umriss x 2.37–21.63, y 1.99–22.01 px; 4 Kontur(en)
\newcommand\icongear{\fill[nonzero rule]
    (3.352,7.008) .. controls (3.133,7.367) and (2.941,7.742) .. (2.777,8.133) .. controls (2.613,8.523) and (2.477,8.922) .. (2.367,9.328) .. controls (2.617,9.453) and (2.844,9.609) .. (3.047,9.797) .. controls (3.25,9.984) and (3.422,10.195) .. (3.562,10.43) .. controls (3.703,10.664) and (3.812,10.914) .. (3.891,11.18) .. controls (3.969,11.445) and (4.008,11.719) .. (4.008,12) .. controls (4.008,12.281) and (3.969,12.555) .. (3.891,12.82) .. controls (3.812,13.086) and (3.703,13.336) .. (3.562,13.57) .. controls (3.422,13.805) and (3.25,14.016) .. (3.047,14.203) .. controls (2.844,14.391) and (2.617,14.547) .. (2.367,14.672) .. controls (2.586,15.484) and (2.91,16.258) .. (3.34,16.992) .. controls (3.77,17.727) and (4.281,18.398) .. (4.875,19.008) .. controls (5.094,18.852) and (5.336,18.734) .. (5.602,18.656) .. controls (5.867,18.578) and (6.141,18.531) .. (6.422,18.516) .. controls (6.703,18.516) and (6.977,18.551) .. (7.242,18.621) .. controls (7.508,18.691) and (7.758,18.797) .. (7.992,18.938) .. controls (8.242,19.062) and (8.461,19.227) .. (8.648,19.43) .. controls (8.836,19.633) and (9,19.852) .. (9.141,20.086) .. controls (9.266,20.336) and (9.359,20.594) .. (9.422,20.859) .. controls (9.484,21.125) and (9.508,21.398) .. (9.492,21.68) .. controls (10.32,21.898) and (11.156,22.008) .. (12,22.008) .. controls (12.844,22.008) and (13.68,21.898) .. (14.508,21.68) .. controls (14.492,21.398) and (14.516,21.125) .. (14.578,20.859) .. controls (14.641,20.594) and (14.734,20.336) .. (14.859,20.086) .. controls (15,19.852) and (15.164,19.633) .. (15.352,19.43) .. controls (15.539,19.227) and (15.758,19.062) .. (16.008,18.938) .. controls (16.242,18.797) and (16.492,18.691) .. (16.758,18.621) .. controls (17.023,18.551) and (17.297,18.516) .. (17.578,18.516) .. controls (17.859,18.531) and (18.133,18.578) .. (18.398,18.656) .. controls (18.664,18.734) and (18.906,18.852) .. (19.125,19.008) .. controls (19.422,18.711) and (19.699,18.395) .. (19.957,18.059) .. controls (20.215,17.723) and (20.445,17.367) .. (20.648,16.992) .. controls (20.867,16.617) and (21.059,16.238) .. (21.223,15.855) .. controls (21.387,15.473) and (21.523,15.078) .. (21.633,14.672) .. controls (21.383,14.547) and (21.156,14.391) .. (20.953,14.203) .. controls (20.75,14.016) and (20.578,13.805) .. (20.438,13.57) .. controls (20.297,13.336) and (20.188,13.086) .. (20.109,12.82) .. controls (20.031,12.555) and (19.992,12.281) .. (19.992,12) .. controls (19.992,11.719) and (20.031,11.445) .. (20.109,11.18) .. controls (20.188,10.914) and (20.297,10.664) .. (20.438,10.43) .. controls (20.578,10.195) and (20.75,9.984) .. (20.953,9.797) .. controls (21.156,9.609) and (21.383,9.453) .. (21.633,9.328) .. controls (21.414,8.516) and (21.09,7.742) .. (20.66,7.008) .. controls (20.23,6.273) and (19.719,5.602) .. (19.125,4.992) .. controls (18.906,5.148) and (18.664,5.266) .. (18.398,5.344) .. controls (18.133,5.422) and (17.859,5.469) .. (17.578,5.484) .. controls (17.297,5.484) and (17.023,5.449) .. (16.758,5.379) .. controls (16.492,5.309) and (16.242,5.203) .. (16.008,5.062) .. controls (15.758,4.938) and (15.539,4.773) .. (15.352,4.57) .. controls (15.164,4.367) and (15,4.148) .. (14.859,3.914) .. controls (14.734,3.664) and (14.641,3.406) .. (14.578,3.141) .. controls (14.516,2.875) and (14.492,2.602) .. (14.508,2.32) .. controls (13.68,2.102) and (12.844,1.992) .. (12,1.992) .. controls (11.156,1.992) and (10.32,2.102) .. (9.492,2.32) .. controls (9.508,2.602) and (9.484,2.875) .. (9.422,3.141) .. controls (9.359,3.406) and (9.266,3.664) .. (9.141,3.914) .. controls (9,4.148) and (8.836,4.367) .. (8.648,4.57) .. controls (8.461,4.773) and (8.242,4.938) .. (7.992,5.062) .. controls (7.758,5.203) and (7.508,5.309) .. (7.242,5.379) .. controls (6.977,5.449) and (6.703,5.484) .. (6.422,5.484) .. controls (6.141,5.469) and (5.867,5.422) .. (5.602,5.344) .. controls (5.336,5.266) and (5.094,5.148) .. (4.875,4.992) .. controls (4.578,5.289) and (4.301,5.605) .. (4.043,5.941) .. controls (3.785,6.277) and (3.555,6.633) .. (3.352,7.008) -- cycle
    (9,6.797) .. controls (9.531,6.5) and (9.992,6.113) .. (10.383,5.637) .. controls (10.773,5.16) and (11.062,4.625) .. (11.25,4.031) .. controls (11.5,4.016) and (11.75,4.008) .. (12,4.008) .. controls (12.25,4.008) and (12.5,4.016) .. (12.75,4.031) .. controls (12.938,4.609) and (13.227,5.141) .. (13.617,5.625) .. controls (14.008,6.109) and (14.469,6.5) .. (15,6.797) .. controls (15.531,7.109) and (16.102,7.312) .. (16.711,7.406) .. controls (17.32,7.5) and (17.922,7.484) .. (18.516,7.359) .. controls (18.672,7.562) and (18.812,7.773) .. (18.938,7.992) .. controls (19.062,8.211) and (19.172,8.438) .. (19.266,8.672) .. controls (18.859,9.125) and (18.547,9.641) .. (18.328,10.219) .. controls (18.109,10.797) and (18,11.391) .. (18,12) .. controls (18,12.625) and (18.113,13.227) .. (18.34,13.805) .. controls (18.566,14.383) and (18.875,14.891) .. (19.266,15.328) .. controls (19.172,15.562) and (19.062,15.789) .. (18.938,16.008) .. controls (18.812,16.227) and (18.672,16.438) .. (18.516,16.641) .. controls (17.922,16.516) and (17.32,16.5) .. (16.711,16.594) .. controls (16.102,16.688) and (15.531,16.891) .. (15,17.203) .. controls (14.469,17.5) and (14.008,17.887) .. (13.617,18.363) .. controls (13.227,18.84) and (12.938,19.375) .. (12.75,19.969) .. controls (12.5,19.984) and (12.25,19.992) .. (12,19.992) .. controls (11.75,19.992) and (11.5,19.984) .. (11.25,19.969) .. controls (11.062,19.391) and (10.773,18.859) .. (10.383,18.375) .. controls (9.992,17.891) and (9.531,17.5) .. (9,17.203) .. controls (8.469,16.891) and (7.898,16.688) .. (7.289,16.594) .. controls (6.68,16.5) and (6.078,16.516) .. (5.484,16.641) .. controls (5.328,16.438) and (5.188,16.227) .. (5.062,16.008) .. controls (4.938,15.789) and (4.828,15.562) .. (4.734,15.328) .. controls (5.141,14.875) and (5.453,14.359) .. (5.672,13.781) .. controls (5.891,13.203) and (6,12.609) .. (6,12) .. controls (6,11.375) and (5.887,10.773) .. (5.66,10.195) .. controls (5.434,9.617) and (5.125,9.109) .. (4.734,8.672) .. controls (4.828,8.438) and (4.938,8.211) .. (5.062,7.992) .. controls (5.188,7.773) and (5.328,7.562) .. (5.484,7.359) .. controls (6.078,7.484) and (6.68,7.5) .. (7.289,7.406) .. controls (7.898,7.312) and (8.469,7.109) .. (9,6.797) -- cycle
    (12,9) .. controls (11.609,9) and (11.227,9.074) .. (10.852,9.223) .. controls (10.477,9.371) and (10.148,9.586) .. (9.867,9.867) .. controls (9.586,10.148) and (9.371,10.477) .. (9.223,10.852) .. controls (9.074,11.227) and (9,11.609) .. (9,12) .. controls (9,12.391) and (9.074,12.773) .. (9.223,13.148) .. controls (9.371,13.523) and (9.586,13.852) .. (9.867,14.133) .. controls (10.148,14.414) and (10.477,14.629) .. (10.852,14.777) .. controls (11.227,14.926) and (11.609,15) .. (12,15) .. controls (12.391,15) and (12.773,14.926) .. (13.148,14.777) .. controls (13.523,14.629) and (13.852,14.414) .. (14.133,14.133) .. controls (14.414,13.852) and (14.629,13.523) .. (14.777,13.148) .. controls (14.926,12.773) and (15,12.391) .. (15,12) .. controls (15,11.609) and (14.926,11.227) .. (14.777,10.852) .. controls (14.629,10.477) and (14.414,10.148) .. (14.133,9.867) .. controls (13.852,9.586) and (13.523,9.371) .. (13.148,9.223) .. controls (12.773,9.074) and (12.391,9) .. (12,9) -- cycle
    (12,10.992) .. controls (12.125,10.992) and (12.25,11.02) .. (12.375,11.074) .. controls (12.5,11.129) and (12.609,11.203) .. (12.703,11.297) .. controls (12.797,11.391) and (12.871,11.5) .. (12.926,11.625) .. controls (12.98,11.75) and (13.008,11.875) .. (13.008,12) .. controls (13.008,12.125) and (12.98,12.25) .. (12.926,12.375) .. controls (12.871,12.5) and (12.797,12.609) .. (12.703,12.703) .. controls (12.609,12.797) and (12.5,12.871) .. (12.375,12.926) .. controls (12.25,12.98) and (12.125,13.008) .. (12,13.008) .. controls (11.875,13.008) and (11.75,12.98) .. (11.625,12.926) .. controls (11.5,12.871) and (11.391,12.797) .. (11.297,12.703) .. controls (11.203,12.609) and (11.129,12.5) .. (11.074,12.375) .. controls (11.02,12.25) and (10.992,12.125) .. (10.992,12) .. controls (10.992,11.875) and (11.02,11.75) .. (11.074,11.625) .. controls (11.129,11.5) and (11.203,11.391) .. (11.297,11.297) .. controls (11.391,11.203) and (11.5,11.129) .. (11.625,11.074) .. controls (11.75,11.02) and (11.875,10.992) .. (12,10.992) -- cycle;}

% U+E90C uniE90C – icon-orders (Liste mit Plus); Vorschub 24.000 px; Umriss x 3.00–21.00, y 1.01–19.99 px; 4 Kontur(en)
\newcommand\iconorders{\fill[nonzero rule]
    (18,9) -- (18,6) -- (21,6) -- (21,4.008) -- (18,4.008) -- (18,1.008) -- (16.008,1.008) -- (16.008,4.008) -- (13.008,4.008) -- (13.008,6) -- (16.008,6) -- (16.008,9) -- cycle
    (10.992,6) -- (10.992,4.008) -- (3,4.008) -- (3,6) -- cycle
    (21,13.008) -- (21,10.992) -- (3,10.992) -- (3,13.008) -- (21,13.008) -- cycle
    (21,19.992) -- (21,18) -- (3,18) -- (3,19.992) -- cycle;}

% U+E93A uniE93A – icon-information (Kreis mit i); Vorschub 24.000 px; Umriss x 1.99–22.01, y 1.99–22.01 px; 4 Kontur(en)
\newcommand\iconinfo{\fill[nonzero rule]
    (12,1.992) .. controls (10.625,1.992) and (9.328,2.258) .. (8.109,2.789) .. controls (6.891,3.305) and (5.828,4.016) .. (4.922,4.922) .. controls (4.016,5.828) and (3.305,6.891) .. (2.789,8.109) .. controls (2.258,9.328) and (1.992,10.625) .. (1.992,12) .. controls (1.992,13.375) and (2.258,14.672) .. (2.789,15.891) .. controls (3.305,17.109) and (4.016,18.172) .. (4.922,19.078) .. controls (5.828,19.984) and (6.891,20.695) .. (8.109,21.211) .. controls (9.328,21.742) and (10.625,22.008) .. (12,22.008) .. controls (13.375,22.008) and (14.672,21.742) .. (15.891,21.211) .. controls (17.109,20.695) and (18.172,19.984) .. (19.078,19.078) .. controls (19.984,18.172) and (20.695,17.109) .. (21.211,15.891) .. controls (21.742,14.672) and (22.008,13.375) .. (22.008,12) .. controls (22.008,10.625) and (21.742,9.328) .. (21.211,8.109) .. controls (20.695,6.891) and (19.984,5.828) .. (19.078,4.922) .. controls (18.172,4.016) and (17.109,3.305) .. (15.891,2.789) .. controls (14.672,2.258) and (13.375,1.992) .. (12,1.992) -- cycle
    (12,4.008) .. controls (13.062,4.008) and (14.082,4.211) .. (15.059,4.617) .. controls (16.035,5.023) and (16.898,5.602) .. (17.648,6.352) .. controls (18.398,7.102) and (18.977,7.965) .. (19.383,8.941) .. controls (19.789,9.918) and (19.992,10.938) .. (19.992,12) .. controls (19.992,13.062) and (19.789,14.082) .. (19.383,15.059) .. controls (18.977,16.035) and (18.398,16.898) .. (17.648,17.648) .. controls (16.898,18.398) and (16.035,18.977) .. (15.059,19.383) .. controls (14.082,19.789) and (13.062,19.992) .. (12,19.992) .. controls (10.938,19.992) and (9.918,19.789) .. (8.941,19.383) .. controls (7.965,18.977) and (7.102,18.398) .. (6.352,17.648) .. controls (5.602,16.898) and (5.023,16.035) .. (4.617,15.059) .. controls (4.211,14.082) and (4.008,13.062) .. (4.008,12) .. controls (4.008,10.938) and (4.211,9.918) .. (4.617,8.941) .. controls (5.023,7.965) and (5.602,7.102) .. (6.352,6.352) .. controls (7.102,5.602) and (7.965,5.023) .. (8.941,4.617) .. controls (9.918,4.211) and (10.938,4.008) .. (12,4.008) -- cycle
    (10.992,16.992) -- (13.008,16.992) -- (13.008,15) -- (10.992,15) -- (10.992,16.992) -- cycle
    (10.992,13.008) -- (13.008,13.008) -- (13.008,7.008) -- (10.992,7.008) -- (10.992,13.008) -- cycle;}


\newif\ifmarken \markentrue
\ifdefined\ohnemarken \markenfalse\fi

\newcommand\breite{1484}% 1704 - 70 (Bildspalte) - 150 (Zahnrad, Schließen)
\newcommand\hoehe{145}%   zweite Zeile endet bei 135, + 10 px
\newcommand\links{-70}%   Verschiebung der Tabellenspalten (ID beginnt bei x10)
\newcommand\rand{1474}%   rechter Tabellenrand (\breite - 10)

% Textbreiten in px (vor dem Bild messen; in TikZ ist \nullfont aktiv)
\newlength\tw
\makeatletter
\newcommand\breitepx[2]{\settowidth\tw{\schrift #2}%
  \edef#1{\strip@pt\dimexpr\tw*4/3*7200/7227\relax}}
\newcommand\probe[2]{\breitepx\probewert{#1}%
  \typeout{POSITIONEN-VOLL: '#1' \probewert\space px (gemessen #2)}}
\makeatother
\probe{Active (0)}{55,0}\probe{Pending (0)}{66,1}\probe{History}{40,6}
\probe{ID}{11,9}\probe{Instrument}{60,7}\probe{Type}{28,1}\probe{Amount}{44,1}
\probe{Trade Volume}{78,4}\probe{Open Rate}{59,2}\probe{Open Time}{62,2}
\probe{S/L}{18,8}\probe{T/P}{19,1}\probe{Swap}{31,6}
\probe{Commission}{69,8}\probe{Total Profit}{61,5}\probe{Margin}{39,3}
\probe{All sides}{47,4}\probe{All profits}{53,8}
\breitepx\probewert{Search}\typeout{POSITIONEN-VOLL: 'Search' \probewert\space px (nicht gemessen)}
\breitepx\wn{\halbfett New order}\pgfmathsetmacro\probewert{97.7-\wn-22.3-12}
\typeout{POSITIONEN-VOLL: 'New order' (600) \wn\space px; Abstand Plus-Text \probewert\space px}
\breitepx\wa{Active}\breitepx\wp{Pending}\breitepx\wh{History}
\pgfmathsetmacro\xh{\rand-(\wh+32)}% Reiter rechtsbündig am Tabellenrand
\pgfmathsetmacro\xp{\xh-(\wp+32)}
\pgfmathsetmacro\xa{\xp-(\wa+32)}
\typeout{POSITIONEN-VOLL: Reiter ab x \xa, \xp, \xh\space bis \rand\space px}

% \flaeche{<Zeilenoberkante>}: getönte Zeile, Zellen-Fläche unter dem 1,6-px-Rand;
% außen Radius 3 px (oben innen 3 x 1,4 px)
\newcommand\flaeche[1]{%
  \fill[aktiv,fill opacity=0.1] (\breite/2,#1+1.6) -- (13,#1+1.6)
    arc[start angle=270,end angle=180,x radius=3,y radius=1.4] -- (10,#1+26.6)
    arc[start angle=180,end angle=90,radius=3] -- (\rand-3,#1+29.6)
    arc[start angle=90,end angle=0,radius=3] -- (\rand,#1+3)
    arc[start angle=0,end angle=-90,x radius=3,y radius=1.4] -- cycle;}

% \kopf{<Original-x des span>}{<Text>} / \zelle{<Mitte>}{<y>}{<Farbe>}{<Text>}
\newcommand\kopf[2]{\node[anchor=base west,text=label] at (#1,65.4) {#2};}
\newcommand\zelle[4]{\node[anchor=base,text=#3] at (#1,#2) {#4};}

% \icon{<x>}{<y Oberkante>}{<Größe px>}{<Farbe>}{<Icon>}: Box Größe x Größe
% (icons.tikz: 24 px, Ursprung unten links, y nach oben)
\newcommand\icon[5]{%
  \begin{scope}[shift={(#1,#2+#3)},x=#3/32*1bp,y=#3/32*1bp,color=#4]#5\end{scope}}

% \marke{<Original-x des Spaltennamens>}{<Nummer>}: Kreis 17 px, Mitte 3 px + Radius davor
\newcommand\marke[2]{%
  \fill[white] (#1-11.5,61) circle[radius=8.5];
  \node[anchor=base,text=markenziffer,font=\fontsize{6.954bp}{8bp}\halbfett]
    at (#1-11.5,61+3.091) {#2};}% 12 px * 17/22

\begin{document}%
\begin{tikzpicture}[x=\px,y={(0,-\px)},font=\schrift,
    every node/.style={inner sep=0,outer sep=0}]
  \path[use as bounding box] (0,0) rectangle (\breite,\hoehe);
  % Block und Zeilen
  \fill[block,rounded corners=7\px] (0,0) rectangle (\breite,\hoehe);
  \flaeche{75.8}\flaeche{105.4}
  % Leiste oben links (Lage wie gemessen)
  \draw[gewinn,line width=0.8\px,rounded corners=4.6\px]
    (10.4,10.6) rectangle (109.6,39.4);% Rand 0,8 px innen, Radius 5 px außen
  \ifdefined\iconplus \icon{22.3}{17.8}{12}{white}{\iconplus}\fi
  \node[anchor=base east,text=white,font=\schrift\halbfett] at (97.7,29.0) {New order};
  \fill[reiter,rounded corners=4\px] (120,10) rectangle (265,40);
  \node[anchor=base west,text=suchtext] at (130.8,29.2) {Search};
  \ifdefined\iconsearch \icon{238}{17.5}{15}{label}{\iconsearch}\fi
  \fill[reiter,rounded corners=4\px] (275,10) rectangle (420,40);
  \node[anchor=base west,text=white] at (285.8,29.4) {All sides};
  \fill[reiter,rounded corners=4\px] (430,10) rectangle (575,40);
  \node[anchor=base west,text=white] at (440.8,29.4) {All profits};
  % Reiter (ohne Zähler), bündig mit dem rechten Tabellenrand
  \fill[reiter,rounded corners=3\px] (\xa,10) rectangle (\rand,40);
  \fill[aktiv] (\xp,10) {[rounded corners=3\px] -- (\xa,10) -- (\xa,40)} -- (\xp,40) -- cycle;
  \node[anchor=base,text=white] at ({\xa+(\wa+32)/2},29.2) {Active};
  \node[anchor=base,text=white] at ({\xp+(\wp+32)/2},29.2) {Pending};
  \node[anchor=base,text=white] at ({\xh+(\wh+32)/2},29.2) {History};
  % Tabelle (Original-x, um die Bildspalte nach links verschoben)
  \begin{scope}[xshift=\links\px]
    \kopf{102.5}{ID}\kopf{177}{Instrument}\kopf{298}{Type}\kopf{383.8}{Amount}
    \kopf{493.5}{Trade Volume}\kopf{640}{Open Rate}\kopf{787.5}{Open Time}
    \kopf{941}{S/L}\kopf{1030.9}{T/P}\kopf{1113.6}{Swap}\kopf{1207.4}{Commission}
    \kopf{1344.2}{Total Profit}\kopf{1467.9}{Margin}
    \foreach \y/\farbe in {95.0/gewinn,124.6/verlust} {
      \zelle{115.9}{\y}{white}{Nummer}
      \zelle{209.4}{\y}{white}{Instrument}% Mitte 214,9 - 5,5 (Pfeil-Bild weggelassen)
      \zelle{319.6}{\y}{white}{Buy}
      \zelle{413.35}{\y}{white}{Menge}
      \zelle{540.25}{\y}{white}{Wert}
      \zelle{677.1}{\y}{white}{Kurs}
      \zelle{826.1}{\y}{white}{Zeitpunkt}
      \zelle{957.9}{\y}{white}{—}
      \zelle{1047.9}{\y}{white}{—}
      \zelle{1136.9}{\y}{white}{Betrag}
      \zelle{1249.85}{\y}{white}{—}
      \zelle{1382.5}{\y}{\farbe}{Ergebnis}
      \zelle{1495}{\y}{white}{Betrag}}
    \ifmarken
      \marke{1113.6}{1}
      \marke{1344.2}{2}
      \marke{1467.9}{3}
    \fi
  \end{scope}
\end{tikzpicture}%
\end{document}
'):
% Stil wie info-desktop.tex (weißer Kreis 17 px, Ziffer Inter 600 in #04273C), mittig auf
% der Kopfzeile, 3 px vor dem Spaltennamen: 1 Swap, 2 Total Profit, 3 Margin.
\documentclass[border=0pt]{standalone}
\usepackage{fontspec}
\usepackage{tikz}
\setmainfont{Inter}% Inter 4.0 (Paket fonts-inter), wie bar.tex
\newfontfamily\halbfett{Inter SemiBold}% New order, Ziffern der Marken (wie panel.pdf)

\newlength\px \setlength\px{0.75bp}% 1 CSS-px
\newcommand\schrift{\fontsize{9bp}{10.5bp}\selectfont}% 12 px / line-height 14 px
\definecolor{block}{RGB}{4,39,60}%      .maintable
\definecolor{reiter}{RGB}{0,19,40}%     ul.nav-tabs, input, ng-select-container
\definecolor{aktiv}{RGB}{14,80,114}%    a.nav-link.active; Zeilen mit 10 % Deckung
\definecolor{label}{RGB}{168,187,190}%  th, i.icon-search
\definecolor{gewinn}{RGB}{83,188,81}%   color-green-main, Rand von New order
\definecolor{verlust}{RGB}{249,64,86}%  color-red-main
\definecolor{suchtext}{RGB}{168,187,190}% Platzhalter „Search“ – NICHT GEMESSEN (vorläufig)
\definecolor{markenziffer}{HTML}{04273C}% Ziffer der Marken (wie panel.pdf)

% Icons als TikZ-Pfade – erzeugt von tools/icons_aus_font.py, nicht von Hand ändern.
% Schrift: 230d1684-icomoon-1-icomoon.b9abfeda36b81d3e.ttf (icomoon, Version 1.0, TrueType), sha256 951e3a5eca89112f…
% Herkunft: Webfont „icomoon“ der Plattform (icomoon.b9abfeda36b81d3e.ttf; icomoon.47a83e5a26e26197.woff hat dieselben Tabellen und Glyphen), vom Nutzer hochgeladen am 02.10.2026. Die Schriftdateien liegen NICHT im Repo (Lizenz unklar), nur diese Pfade.
% unitsPerEm 1024; Metriken hhea: Ascent 960, Descent -64.
% Box wie das <i>-Element im Browser: font-size 24px, line-height 24px
% -> Box 24 x 24 px, Ursprung unten links, y nach oben;
%    Grundlinie 1.500 px über der Boxunterkante. 1 Einheit = 1 CSS-px.
% Füllregel nonzero (wie im Browser). Farbe kommt von außen (color=…).
% Gebrauch: \begin{scope}[shift={(<x>,<y>)},x=\px,y=\px,color=<Farbe>]\iconbell\end{scope}

% U+E938 uniE938 – icon-notifications (Glocke); Vorschub 24.000 px; Umriss x 1.99–22.01, y 1.99–22.01 px; 3 Kontur(en)
\newcommand\iconbell{\fill[nonzero rule]
    (19.992,7.711) -- (22.008,7.711) -- (22.008,5.812) -- (1.992,5.812) -- (1.992,7.711) -- (4.008,7.711) -- (4.008,14.391) .. controls (4.008,15.391) and (4.211,16.359) .. (4.617,17.297) .. controls (5.023,18.234) and (5.602,19.055) .. (6.352,19.758) .. controls (7.102,20.477) and (7.965,21.031) .. (8.941,21.422) .. controls (9.918,21.812) and (10.938,22.008) .. (12,22.008) .. controls (13.062,22.008) and (14.082,21.812) .. (15.059,21.422) .. controls (16.035,21.031) and (16.898,20.477) .. (17.648,19.758) .. controls (18.398,19.055) and (18.977,18.234) .. (19.383,17.297) .. controls (19.789,16.359) and (19.992,15.391) .. (19.992,14.391) -- (19.992,7.711) -- cycle
    (18,7.711) -- (18,14.391) .. controls (18,15.141) and (17.848,15.867) .. (17.543,16.57) .. controls (17.238,17.273) and (16.805,17.891) .. (16.242,18.422) .. controls (15.68,18.953) and (15.031,19.363) .. (14.297,19.652) .. controls (13.562,19.941) and (12.797,20.086) .. (12,20.086) .. controls (11.203,20.086) and (10.438,19.941) .. (9.703,19.652) .. controls (8.969,19.363) and (8.32,18.953) .. (7.758,18.422) .. controls (7.195,17.891) and (6.762,17.273) .. (6.457,16.57) .. controls (6.152,15.867) and (6,15.141) .. (6,14.391) -- (6,7.711) -- cycle
    (9,3.891) -- (15,3.891) -- (15,1.992) -- (9,1.992) -- cycle;}

% U+E937 uniE937 – icon-settings (Zahnrad); Vorschub 24.000 px; Umriss x 2.37–21.63, y 1.99–22.01 px; 4 Kontur(en)
\newcommand\icongear{\fill[nonzero rule]
    (3.352,7.008) .. controls (3.133,7.367) and (2.941,7.742) .. (2.777,8.133) .. controls (2.613,8.523) and (2.477,8.922) .. (2.367,9.328) .. controls (2.617,9.453) and (2.844,9.609) .. (3.047,9.797) .. controls (3.25,9.984) and (3.422,10.195) .. (3.562,10.43) .. controls (3.703,10.664) and (3.812,10.914) .. (3.891,11.18) .. controls (3.969,11.445) and (4.008,11.719) .. (4.008,12) .. controls (4.008,12.281) and (3.969,12.555) .. (3.891,12.82) .. controls (3.812,13.086) and (3.703,13.336) .. (3.562,13.57) .. controls (3.422,13.805) and (3.25,14.016) .. (3.047,14.203) .. controls (2.844,14.391) and (2.617,14.547) .. (2.367,14.672) .. controls (2.586,15.484) and (2.91,16.258) .. (3.34,16.992) .. controls (3.77,17.727) and (4.281,18.398) .. (4.875,19.008) .. controls (5.094,18.852) and (5.336,18.734) .. (5.602,18.656) .. controls (5.867,18.578) and (6.141,18.531) .. (6.422,18.516) .. controls (6.703,18.516) and (6.977,18.551) .. (7.242,18.621) .. controls (7.508,18.691) and (7.758,18.797) .. (7.992,18.938) .. controls (8.242,19.062) and (8.461,19.227) .. (8.648,19.43) .. controls (8.836,19.633) and (9,19.852) .. (9.141,20.086) .. controls (9.266,20.336) and (9.359,20.594) .. (9.422,20.859) .. controls (9.484,21.125) and (9.508,21.398) .. (9.492,21.68) .. controls (10.32,21.898) and (11.156,22.008) .. (12,22.008) .. controls (12.844,22.008) and (13.68,21.898) .. (14.508,21.68) .. controls (14.492,21.398) and (14.516,21.125) .. (14.578,20.859) .. controls (14.641,20.594) and (14.734,20.336) .. (14.859,20.086) .. controls (15,19.852) and (15.164,19.633) .. (15.352,19.43) .. controls (15.539,19.227) and (15.758,19.062) .. (16.008,18.938) .. controls (16.242,18.797) and (16.492,18.691) .. (16.758,18.621) .. controls (17.023,18.551) and (17.297,18.516) .. (17.578,18.516) .. controls (17.859,18.531) and (18.133,18.578) .. (18.398,18.656) .. controls (18.664,18.734) and (18.906,18.852) .. (19.125,19.008) .. controls (19.422,18.711) and (19.699,18.395) .. (19.957,18.059) .. controls (20.215,17.723) and (20.445,17.367) .. (20.648,16.992) .. controls (20.867,16.617) and (21.059,16.238) .. (21.223,15.855) .. controls (21.387,15.473) and (21.523,15.078) .. (21.633,14.672) .. controls (21.383,14.547) and (21.156,14.391) .. (20.953,14.203) .. controls (20.75,14.016) and (20.578,13.805) .. (20.438,13.57) .. controls (20.297,13.336) and (20.188,13.086) .. (20.109,12.82) .. controls (20.031,12.555) and (19.992,12.281) .. (19.992,12) .. controls (19.992,11.719) and (20.031,11.445) .. (20.109,11.18) .. controls (20.188,10.914) and (20.297,10.664) .. (20.438,10.43) .. controls (20.578,10.195) and (20.75,9.984) .. (20.953,9.797) .. controls (21.156,9.609) and (21.383,9.453) .. (21.633,9.328) .. controls (21.414,8.516) and (21.09,7.742) .. (20.66,7.008) .. controls (20.23,6.273) and (19.719,5.602) .. (19.125,4.992) .. controls (18.906,5.148) and (18.664,5.266) .. (18.398,5.344) .. controls (18.133,5.422) and (17.859,5.469) .. (17.578,5.484) .. controls (17.297,5.484) and (17.023,5.449) .. (16.758,5.379) .. controls (16.492,5.309) and (16.242,5.203) .. (16.008,5.062) .. controls (15.758,4.938) and (15.539,4.773) .. (15.352,4.57) .. controls (15.164,4.367) and (15,4.148) .. (14.859,3.914) .. controls (14.734,3.664) and (14.641,3.406) .. (14.578,3.141) .. controls (14.516,2.875) and (14.492,2.602) .. (14.508,2.32) .. controls (13.68,2.102) and (12.844,1.992) .. (12,1.992) .. controls (11.156,1.992) and (10.32,2.102) .. (9.492,2.32) .. controls (9.508,2.602) and (9.484,2.875) .. (9.422,3.141) .. controls (9.359,3.406) and (9.266,3.664) .. (9.141,3.914) .. controls (9,4.148) and (8.836,4.367) .. (8.648,4.57) .. controls (8.461,4.773) and (8.242,4.938) .. (7.992,5.062) .. controls (7.758,5.203) and (7.508,5.309) .. (7.242,5.379) .. controls (6.977,5.449) and (6.703,5.484) .. (6.422,5.484) .. controls (6.141,5.469) and (5.867,5.422) .. (5.602,5.344) .. controls (5.336,5.266) and (5.094,5.148) .. (4.875,4.992) .. controls (4.578,5.289) and (4.301,5.605) .. (4.043,5.941) .. controls (3.785,6.277) and (3.555,6.633) .. (3.352,7.008) -- cycle
    (9,6.797) .. controls (9.531,6.5) and (9.992,6.113) .. (10.383,5.637) .. controls (10.773,5.16) and (11.062,4.625) .. (11.25,4.031) .. controls (11.5,4.016) and (11.75,4.008) .. (12,4.008) .. controls (12.25,4.008) and (12.5,4.016) .. (12.75,4.031) .. controls (12.938,4.609) and (13.227,5.141) .. (13.617,5.625) .. controls (14.008,6.109) and (14.469,6.5) .. (15,6.797) .. controls (15.531,7.109) and (16.102,7.312) .. (16.711,7.406) .. controls (17.32,7.5) and (17.922,7.484) .. (18.516,7.359) .. controls (18.672,7.562) and (18.812,7.773) .. (18.938,7.992) .. controls (19.062,8.211) and (19.172,8.438) .. (19.266,8.672) .. controls (18.859,9.125) and (18.547,9.641) .. (18.328,10.219) .. controls (18.109,10.797) and (18,11.391) .. (18,12) .. controls (18,12.625) and (18.113,13.227) .. (18.34,13.805) .. controls (18.566,14.383) and (18.875,14.891) .. (19.266,15.328) .. controls (19.172,15.562) and (19.062,15.789) .. (18.938,16.008) .. controls (18.812,16.227) and (18.672,16.438) .. (18.516,16.641) .. controls (17.922,16.516) and (17.32,16.5) .. (16.711,16.594) .. controls (16.102,16.688) and (15.531,16.891) .. (15,17.203) .. controls (14.469,17.5) and (14.008,17.887) .. (13.617,18.363) .. controls (13.227,18.84) and (12.938,19.375) .. (12.75,19.969) .. controls (12.5,19.984) and (12.25,19.992) .. (12,19.992) .. controls (11.75,19.992) and (11.5,19.984) .. (11.25,19.969) .. controls (11.062,19.391) and (10.773,18.859) .. (10.383,18.375) .. controls (9.992,17.891) and (9.531,17.5) .. (9,17.203) .. controls (8.469,16.891) and (7.898,16.688) .. (7.289,16.594) .. controls (6.68,16.5) and (6.078,16.516) .. (5.484,16.641) .. controls (5.328,16.438) and (5.188,16.227) .. (5.062,16.008) .. controls (4.938,15.789) and (4.828,15.562) .. (4.734,15.328) .. controls (5.141,14.875) and (5.453,14.359) .. (5.672,13.781) .. controls (5.891,13.203) and (6,12.609) .. (6,12) .. controls (6,11.375) and (5.887,10.773) .. (5.66,10.195) .. controls (5.434,9.617) and (5.125,9.109) .. (4.734,8.672) .. controls (4.828,8.438) and (4.938,8.211) .. (5.062,7.992) .. controls (5.188,7.773) and (5.328,7.562) .. (5.484,7.359) .. controls (6.078,7.484) and (6.68,7.5) .. (7.289,7.406) .. controls (7.898,7.312) and (8.469,7.109) .. (9,6.797) -- cycle
    (12,9) .. controls (11.609,9) and (11.227,9.074) .. (10.852,9.223) .. controls (10.477,9.371) and (10.148,9.586) .. (9.867,9.867) .. controls (9.586,10.148) and (9.371,10.477) .. (9.223,10.852) .. controls (9.074,11.227) and (9,11.609) .. (9,12) .. controls (9,12.391) and (9.074,12.773) .. (9.223,13.148) .. controls (9.371,13.523) and (9.586,13.852) .. (9.867,14.133) .. controls (10.148,14.414) and (10.477,14.629) .. (10.852,14.777) .. controls (11.227,14.926) and (11.609,15) .. (12,15) .. controls (12.391,15) and (12.773,14.926) .. (13.148,14.777) .. controls (13.523,14.629) and (13.852,14.414) .. (14.133,14.133) .. controls (14.414,13.852) and (14.629,13.523) .. (14.777,13.148) .. controls (14.926,12.773) and (15,12.391) .. (15,12) .. controls (15,11.609) and (14.926,11.227) .. (14.777,10.852) .. controls (14.629,10.477) and (14.414,10.148) .. (14.133,9.867) .. controls (13.852,9.586) and (13.523,9.371) .. (13.148,9.223) .. controls (12.773,9.074) and (12.391,9) .. (12,9) -- cycle
    (12,10.992) .. controls (12.125,10.992) and (12.25,11.02) .. (12.375,11.074) .. controls (12.5,11.129) and (12.609,11.203) .. (12.703,11.297) .. controls (12.797,11.391) and (12.871,11.5) .. (12.926,11.625) .. controls (12.98,11.75) and (13.008,11.875) .. (13.008,12) .. controls (13.008,12.125) and (12.98,12.25) .. (12.926,12.375) .. controls (12.871,12.5) and (12.797,12.609) .. (12.703,12.703) .. controls (12.609,12.797) and (12.5,12.871) .. (12.375,12.926) .. controls (12.25,12.98) and (12.125,13.008) .. (12,13.008) .. controls (11.875,13.008) and (11.75,12.98) .. (11.625,12.926) .. controls (11.5,12.871) and (11.391,12.797) .. (11.297,12.703) .. controls (11.203,12.609) and (11.129,12.5) .. (11.074,12.375) .. controls (11.02,12.25) and (10.992,12.125) .. (10.992,12) .. controls (10.992,11.875) and (11.02,11.75) .. (11.074,11.625) .. controls (11.129,11.5) and (11.203,11.391) .. (11.297,11.297) .. controls (11.391,11.203) and (11.5,11.129) .. (11.625,11.074) .. controls (11.75,11.02) and (11.875,10.992) .. (12,10.992) -- cycle;}

% U+E90C uniE90C – icon-orders (Liste mit Plus); Vorschub 24.000 px; Umriss x 3.00–21.00, y 1.01–19.99 px; 4 Kontur(en)
\newcommand\iconorders{\fill[nonzero rule]
    (18,9) -- (18,6) -- (21,6) -- (21,4.008) -- (18,4.008) -- (18,1.008) -- (16.008,1.008) -- (16.008,4.008) -- (13.008,4.008) -- (13.008,6) -- (16.008,6) -- (16.008,9) -- cycle
    (10.992,6) -- (10.992,4.008) -- (3,4.008) -- (3,6) -- cycle
    (21,13.008) -- (21,10.992) -- (3,10.992) -- (3,13.008) -- (21,13.008) -- cycle
    (21,19.992) -- (21,18) -- (3,18) -- (3,19.992) -- cycle;}

% U+E93A uniE93A – icon-information (Kreis mit i); Vorschub 24.000 px; Umriss x 1.99–22.01, y 1.99–22.01 px; 4 Kontur(en)
\newcommand\iconinfo{\fill[nonzero rule]
    (12,1.992) .. controls (10.625,1.992) and (9.328,2.258) .. (8.109,2.789) .. controls (6.891,3.305) and (5.828,4.016) .. (4.922,4.922) .. controls (4.016,5.828) and (3.305,6.891) .. (2.789,8.109) .. controls (2.258,9.328) and (1.992,10.625) .. (1.992,12) .. controls (1.992,13.375) and (2.258,14.672) .. (2.789,15.891) .. controls (3.305,17.109) and (4.016,18.172) .. (4.922,19.078) .. controls (5.828,19.984) and (6.891,20.695) .. (8.109,21.211) .. controls (9.328,21.742) and (10.625,22.008) .. (12,22.008) .. controls (13.375,22.008) and (14.672,21.742) .. (15.891,21.211) .. controls (17.109,20.695) and (18.172,19.984) .. (19.078,19.078) .. controls (19.984,18.172) and (20.695,17.109) .. (21.211,15.891) .. controls (21.742,14.672) and (22.008,13.375) .. (22.008,12) .. controls (22.008,10.625) and (21.742,9.328) .. (21.211,8.109) .. controls (20.695,6.891) and (19.984,5.828) .. (19.078,4.922) .. controls (18.172,4.016) and (17.109,3.305) .. (15.891,2.789) .. controls (14.672,2.258) and (13.375,1.992) .. (12,1.992) -- cycle
    (12,4.008) .. controls (13.062,4.008) and (14.082,4.211) .. (15.059,4.617) .. controls (16.035,5.023) and (16.898,5.602) .. (17.648,6.352) .. controls (18.398,7.102) and (18.977,7.965) .. (19.383,8.941) .. controls (19.789,9.918) and (19.992,10.938) .. (19.992,12) .. controls (19.992,13.062) and (19.789,14.082) .. (19.383,15.059) .. controls (18.977,16.035) and (18.398,16.898) .. (17.648,17.648) .. controls (16.898,18.398) and (16.035,18.977) .. (15.059,19.383) .. controls (14.082,19.789) and (13.062,19.992) .. (12,19.992) .. controls (10.938,19.992) and (9.918,19.789) .. (8.941,19.383) .. controls (7.965,18.977) and (7.102,18.398) .. (6.352,17.648) .. controls (5.602,16.898) and (5.023,16.035) .. (4.617,15.059) .. controls (4.211,14.082) and (4.008,13.062) .. (4.008,12) .. controls (4.008,10.938) and (4.211,9.918) .. (4.617,8.941) .. controls (5.023,7.965) and (5.602,7.102) .. (6.352,6.352) .. controls (7.102,5.602) and (7.965,5.023) .. (8.941,4.617) .. controls (9.918,4.211) and (10.938,4.008) .. (12,4.008) -- cycle
    (10.992,16.992) -- (13.008,16.992) -- (13.008,15) -- (10.992,15) -- (10.992,16.992) -- cycle
    (10.992,13.008) -- (13.008,13.008) -- (13.008,7.008) -- (10.992,7.008) -- (10.992,13.008) -- cycle;}


\newif\ifmarken \markentrue
\ifdefined\ohnemarken \markenfalse\fi

\newcommand\breite{1484}% 1704 - 70 (Bildspalte) - 150 (Zahnrad, Schließen)
\newcommand\hoehe{145}%   zweite Zeile endet bei 135, + 10 px
\newcommand\links{-70}%   Verschiebung der Tabellenspalten (ID beginnt bei x10)
\newcommand\rand{1474}%   rechter Tabellenrand (\breite - 10)

% Textbreiten in px (vor dem Bild messen; in TikZ ist \nullfont aktiv)
\newlength\tw
\makeatletter
\newcommand\breitepx[2]{\settowidth\tw{\schrift #2}%
  \edef#1{\strip@pt\dimexpr\tw*4/3*7200/7227\relax}}
\newcommand\probe[2]{\breitepx\probewert{#1}%
  \typeout{POSITIONEN-VOLL: '#1' \probewert\space px (gemessen #2)}}
\makeatother
\probe{Active (0)}{55,0}\probe{Pending (0)}{66,1}\probe{History}{40,6}
\probe{ID}{11,9}\probe{Instrument}{60,7}\probe{Type}{28,1}\probe{Amount}{44,1}
\probe{Trade Volume}{78,4}\probe{Open Rate}{59,2}\probe{Open Time}{62,2}
\probe{S/L}{18,8}\probe{T/P}{19,1}\probe{Swap}{31,6}
\probe{Commission}{69,8}\probe{Total Profit}{61,5}\probe{Margin}{39,3}
\probe{All sides}{47,4}\probe{All profits}{53,8}
\breitepx\probewert{Search}\typeout{POSITIONEN-VOLL: 'Search' \probewert\space px (nicht gemessen)}
\breitepx\wn{\halbfett New order}\pgfmathsetmacro\probewert{97.7-\wn-22.3-12}
\typeout{POSITIONEN-VOLL: 'New order' (600) \wn\space px; Abstand Plus-Text \probewert\space px}
\breitepx\wa{Active}\breitepx\wp{Pending}\breitepx\wh{History}
\pgfmathsetmacro\xh{\rand-(\wh+32)}% Reiter rechtsbündig am Tabellenrand
\pgfmathsetmacro\xp{\xh-(\wp+32)}
\pgfmathsetmacro\xa{\xp-(\wa+32)}
\typeout{POSITIONEN-VOLL: Reiter ab x \xa, \xp, \xh\space bis \rand\space px}

% \flaeche{<Zeilenoberkante>}: getönte Zeile, Zellen-Fläche unter dem 1,6-px-Rand;
% außen Radius 3 px (oben innen 3 x 1,4 px)
\newcommand\flaeche[1]{%
  \fill[aktiv,fill opacity=0.1] (\breite/2,#1+1.6) -- (13,#1+1.6)
    arc[start angle=270,end angle=180,x radius=3,y radius=1.4] -- (10,#1+26.6)
    arc[start angle=180,end angle=90,radius=3] -- (\rand-3,#1+29.6)
    arc[start angle=90,end angle=0,radius=3] -- (\rand,#1+3)
    arc[start angle=0,end angle=-90,x radius=3,y radius=1.4] -- cycle;}

% \kopf{<Original-x des span>}{<Text>} / \zelle{<Mitte>}{<y>}{<Farbe>}{<Text>}
\newcommand\kopf[2]{\node[anchor=base west,text=label] at (#1,65.4) {#2};}
\newcommand\zelle[4]{\node[anchor=base,text=#3] at (#1,#2) {#4};}

% \icon{<x>}{<y Oberkante>}{<Größe px>}{<Farbe>}{<Icon>}: Box Größe x Größe
% (icons.tikz: 24 px, Ursprung unten links, y nach oben)
\newcommand\icon[5]{%
  \begin{scope}[shift={(#1,#2+#3)},x=#3/32*1bp,y=#3/32*1bp,color=#4]#5\end{scope}}

% \marke{<Original-x des Spaltennamens>}{<Nummer>}: Kreis 17 px, Mitte 3 px + Radius davor
\newcommand\marke[2]{%
  \fill[white] (#1-11.5,61) circle[radius=8.5];
  \node[anchor=base,text=markenziffer,font=\fontsize{6.954bp}{8bp}\halbfett]
    at (#1-11.5,61+3.091) {#2};}% 12 px * 17/22

\begin{document}%
\begin{tikzpicture}[x=\px,y={(0,-\px)},font=\schrift,
    every node/.style={inner sep=0,outer sep=0}]
  \path[use as bounding box] (0,0) rectangle (\breite,\hoehe);
  % Block und Zeilen
  \fill[block,rounded corners=7\px] (0,0) rectangle (\breite,\hoehe);
  \flaeche{75.8}\flaeche{105.4}
  % Leiste oben links (Lage wie gemessen)
  \draw[gewinn,line width=0.8\px,rounded corners=4.6\px]
    (10.4,10.6) rectangle (109.6,39.4);% Rand 0,8 px innen, Radius 5 px außen
  \ifdefined\iconplus \icon{22.3}{17.8}{12}{white}{\iconplus}\fi
  \node[anchor=base east,text=white,font=\schrift\halbfett] at (97.7,29.0) {New order};
  \fill[reiter,rounded corners=4\px] (120,10) rectangle (265,40);
  \node[anchor=base west,text=suchtext] at (130.8,29.2) {Search};
  \ifdefined\iconsearch \icon{238}{17.5}{15}{label}{\iconsearch}\fi
  \fill[reiter,rounded corners=4\px] (275,10) rectangle (420,40);
  \node[anchor=base west,text=white] at (285.8,29.4) {All sides};
  \fill[reiter,rounded corners=4\px] (430,10) rectangle (575,40);
  \node[anchor=base west,text=white] at (440.8,29.4) {All profits};
  % Reiter (ohne Zähler), bündig mit dem rechten Tabellenrand
  \fill[reiter,rounded corners=3\px] (\xa,10) rectangle (\rand,40);
  \fill[aktiv] (\xp,10) {[rounded corners=3\px] -- (\xa,10) -- (\xa,40)} -- (\xp,40) -- cycle;
  \node[anchor=base,text=white] at ({\xa+(\wa+32)/2},29.2) {Active};
  \node[anchor=base,text=white] at ({\xp+(\wp+32)/2},29.2) {Pending};
  \node[anchor=base,text=white] at ({\xh+(\wh+32)/2},29.2) {History};
  % Tabelle (Original-x, um die Bildspalte nach links verschoben)
  \begin{scope}[xshift=\links\px]
    \kopf{102.5}{ID}\kopf{177}{Instrument}\kopf{298}{Type}\kopf{383.8}{Amount}
    \kopf{493.5}{Trade Volume}\kopf{640}{Open Rate}\kopf{787.5}{Open Time}
    \kopf{941}{S/L}\kopf{1030.9}{T/P}\kopf{1113.6}{Swap}\kopf{1207.4}{Commission}
    \kopf{1344.2}{Total Profit}\kopf{1467.9}{Margin}
    \foreach \y/\farbe in {95.0/gewinn,124.6/verlust} {
      \zelle{115.9}{\y}{white}{Nummer}
      \zelle{209.4}{\y}{white}{Instrument}% Mitte 214,9 - 5,5 (Pfeil-Bild weggelassen)
      \zelle{319.6}{\y}{white}{Buy}
      \zelle{413.35}{\y}{white}{Menge}
      \zelle{540.25}{\y}{white}{Wert}
      \zelle{677.1}{\y}{white}{Kurs}
      \zelle{826.1}{\y}{white}{Zeitpunkt}
      \zelle{957.9}{\y}{white}{—}
      \zelle{1047.9}{\y}{white}{—}
      \zelle{1136.9}{\y}{white}{Betrag}
      \zelle{1249.85}{\y}{white}{—}
      \zelle{1382.5}{\y}{\farbe}{Ergebnis}
      \zelle{1495}{\y}{white}{Betrag}}
    \ifmarken
      \marke{1113.6}{1}
      \marke{1344.2}{2}
      \marke{1467.9}{3}
    \fi
  \end{scope}
\end{tikzpicture}%
\end{document}
'):
% Stil wie info-desktop.tex (weißer Kreis 17 px, Ziffer Inter 600 in #04273C), mittig auf
% der Kopfzeile, 3 px vor dem Spaltennamen: 1 Swap, 2 Total Profit, 3 Margin.
\documentclass[border=0pt]{standalone}
\usepackage{fontspec}
\usepackage{tikz}
\setmainfont{Inter}% Inter 4.0 (Paket fonts-inter), wie bar.tex
\newfontfamily\halbfett{Inter SemiBold}% New order, Ziffern der Marken (wie panel.pdf)

\newlength\px \setlength\px{0.75bp}% 1 CSS-px
\newcommand\schrift{\fontsize{9bp}{10.5bp}\selectfont}% 12 px / line-height 14 px
\definecolor{block}{RGB}{4,39,60}%      .maintable
\definecolor{reiter}{RGB}{0,19,40}%     ul.nav-tabs, input, ng-select-container
\definecolor{aktiv}{RGB}{14,80,114}%    a.nav-link.active; Zeilen mit 10 % Deckung
\definecolor{label}{RGB}{168,187,190}%  th, i.icon-search
\definecolor{gewinn}{RGB}{83,188,81}%   color-green-main, Rand von New order
\definecolor{verlust}{RGB}{249,64,86}%  color-red-main
\definecolor{suchtext}{RGB}{168,187,190}% Platzhalter „Search“ – NICHT GEMESSEN (vorläufig)
\definecolor{markenziffer}{HTML}{04273C}% Ziffer der Marken (wie panel.pdf)

% Icons als TikZ-Pfade – erzeugt von tools/icons_aus_font.py, nicht von Hand ändern.
% Schrift: 230d1684-icomoon-1-icomoon.b9abfeda36b81d3e.ttf (icomoon, Version 1.0, TrueType), sha256 951e3a5eca89112f…
% Herkunft: Webfont „icomoon“ der Plattform (icomoon.b9abfeda36b81d3e.ttf; icomoon.47a83e5a26e26197.woff hat dieselben Tabellen und Glyphen), vom Nutzer hochgeladen am 02.10.2026. Die Schriftdateien liegen NICHT im Repo (Lizenz unklar), nur diese Pfade.
% unitsPerEm 1024; Metriken hhea: Ascent 960, Descent -64.
% Box wie das <i>-Element im Browser: font-size 24px, line-height 24px
% -> Box 24 x 24 px, Ursprung unten links, y nach oben;
%    Grundlinie 1.500 px über der Boxunterkante. 1 Einheit = 1 CSS-px.
% Füllregel nonzero (wie im Browser). Farbe kommt von außen (color=…).
% Gebrauch: \begin{scope}[shift={(<x>,<y>)},x=\px,y=\px,color=<Farbe>]\iconbell\end{scope}

% U+E938 uniE938 – icon-notifications (Glocke); Vorschub 24.000 px; Umriss x 1.99–22.01, y 1.99–22.01 px; 3 Kontur(en)
\newcommand\iconbell{\fill[nonzero rule]
    (19.992,7.711) -- (22.008,7.711) -- (22.008,5.812) -- (1.992,5.812) -- (1.992,7.711) -- (4.008,7.711) -- (4.008,14.391) .. controls (4.008,15.391) and (4.211,16.359) .. (4.617,17.297) .. controls (5.023,18.234) and (5.602,19.055) .. (6.352,19.758) .. controls (7.102,20.477) and (7.965,21.031) .. (8.941,21.422) .. controls (9.918,21.812) and (10.938,22.008) .. (12,22.008) .. controls (13.062,22.008) and (14.082,21.812) .. (15.059,21.422) .. controls (16.035,21.031) and (16.898,20.477) .. (17.648,19.758) .. controls (18.398,19.055) and (18.977,18.234) .. (19.383,17.297) .. controls (19.789,16.359) and (19.992,15.391) .. (19.992,14.391) -- (19.992,7.711) -- cycle
    (18,7.711) -- (18,14.391) .. controls (18,15.141) and (17.848,15.867) .. (17.543,16.57) .. controls (17.238,17.273) and (16.805,17.891) .. (16.242,18.422) .. controls (15.68,18.953) and (15.031,19.363) .. (14.297,19.652) .. controls (13.562,19.941) and (12.797,20.086) .. (12,20.086) .. controls (11.203,20.086) and (10.438,19.941) .. (9.703,19.652) .. controls (8.969,19.363) and (8.32,18.953) .. (7.758,18.422) .. controls (7.195,17.891) and (6.762,17.273) .. (6.457,16.57) .. controls (6.152,15.867) and (6,15.141) .. (6,14.391) -- (6,7.711) -- cycle
    (9,3.891) -- (15,3.891) -- (15,1.992) -- (9,1.992) -- cycle;}

% U+E937 uniE937 – icon-settings (Zahnrad); Vorschub 24.000 px; Umriss x 2.37–21.63, y 1.99–22.01 px; 4 Kontur(en)
\newcommand\icongear{\fill[nonzero rule]
    (3.352,7.008) .. controls (3.133,7.367) and (2.941,7.742) .. (2.777,8.133) .. controls (2.613,8.523) and (2.477,8.922) .. (2.367,9.328) .. controls (2.617,9.453) and (2.844,9.609) .. (3.047,9.797) .. controls (3.25,9.984) and (3.422,10.195) .. (3.562,10.43) .. controls (3.703,10.664) and (3.812,10.914) .. (3.891,11.18) .. controls (3.969,11.445) and (4.008,11.719) .. (4.008,12) .. controls (4.008,12.281) and (3.969,12.555) .. (3.891,12.82) .. controls (3.812,13.086) and (3.703,13.336) .. (3.562,13.57) .. controls (3.422,13.805) and (3.25,14.016) .. (3.047,14.203) .. controls (2.844,14.391) and (2.617,14.547) .. (2.367,14.672) .. controls (2.586,15.484) and (2.91,16.258) .. (3.34,16.992) .. controls (3.77,17.727) and (4.281,18.398) .. (4.875,19.008) .. controls (5.094,18.852) and (5.336,18.734) .. (5.602,18.656) .. controls (5.867,18.578) and (6.141,18.531) .. (6.422,18.516) .. controls (6.703,18.516) and (6.977,18.551) .. (7.242,18.621) .. controls (7.508,18.691) and (7.758,18.797) .. (7.992,18.938) .. controls (8.242,19.062) and (8.461,19.227) .. (8.648,19.43) .. controls (8.836,19.633) and (9,19.852) .. (9.141,20.086) .. controls (9.266,20.336) and (9.359,20.594) .. (9.422,20.859) .. controls (9.484,21.125) and (9.508,21.398) .. (9.492,21.68) .. controls (10.32,21.898) and (11.156,22.008) .. (12,22.008) .. controls (12.844,22.008) and (13.68,21.898) .. (14.508,21.68) .. controls (14.492,21.398) and (14.516,21.125) .. (14.578,20.859) .. controls (14.641,20.594) and (14.734,20.336) .. (14.859,20.086) .. controls (15,19.852) and (15.164,19.633) .. (15.352,19.43) .. controls (15.539,19.227) and (15.758,19.062) .. (16.008,18.938) .. controls (16.242,18.797) and (16.492,18.691) .. (16.758,18.621) .. controls (17.023,18.551) and (17.297,18.516) .. (17.578,18.516) .. controls (17.859,18.531) and (18.133,18.578) .. (18.398,18.656) .. controls (18.664,18.734) and (18.906,18.852) .. (19.125,19.008) .. controls (19.422,18.711) and (19.699,18.395) .. (19.957,18.059) .. controls (20.215,17.723) and (20.445,17.367) .. (20.648,16.992) .. controls (20.867,16.617) and (21.059,16.238) .. (21.223,15.855) .. controls (21.387,15.473) and (21.523,15.078) .. (21.633,14.672) .. controls (21.383,14.547) and (21.156,14.391) .. (20.953,14.203) .. controls (20.75,14.016) and (20.578,13.805) .. (20.438,13.57) .. controls (20.297,13.336) and (20.188,13.086) .. (20.109,12.82) .. controls (20.031,12.555) and (19.992,12.281) .. (19.992,12) .. controls (19.992,11.719) and (20.031,11.445) .. (20.109,11.18) .. controls (20.188,10.914) and (20.297,10.664) .. (20.438,10.43) .. controls (20.578,10.195) and (20.75,9.984) .. (20.953,9.797) .. controls (21.156,9.609) and (21.383,9.453) .. (21.633,9.328) .. controls (21.414,8.516) and (21.09,7.742) .. (20.66,7.008) .. controls (20.23,6.273) and (19.719,5.602) .. (19.125,4.992) .. controls (18.906,5.148) and (18.664,5.266) .. (18.398,5.344) .. controls (18.133,5.422) and (17.859,5.469) .. (17.578,5.484) .. controls (17.297,5.484) and (17.023,5.449) .. (16.758,5.379) .. controls (16.492,5.309) and (16.242,5.203) .. (16.008,5.062) .. controls (15.758,4.938) and (15.539,4.773) .. (15.352,4.57) .. controls (15.164,4.367) and (15,4.148) .. (14.859,3.914) .. controls (14.734,3.664) and (14.641,3.406) .. (14.578,3.141) .. controls (14.516,2.875) and (14.492,2.602) .. (14.508,2.32) .. controls (13.68,2.102) and (12.844,1.992) .. (12,1.992) .. controls (11.156,1.992) and (10.32,2.102) .. (9.492,2.32) .. controls (9.508,2.602) and (9.484,2.875) .. (9.422,3.141) .. controls (9.359,3.406) and (9.266,3.664) .. (9.141,3.914) .. controls (9,4.148) and (8.836,4.367) .. (8.648,4.57) .. controls (8.461,4.773) and (8.242,4.938) .. (7.992,5.062) .. controls (7.758,5.203) and (7.508,5.309) .. (7.242,5.379) .. controls (6.977,5.449) and (6.703,5.484) .. (6.422,5.484) .. controls (6.141,5.469) and (5.867,5.422) .. (5.602,5.344) .. controls (5.336,5.266) and (5.094,5.148) .. (4.875,4.992) .. controls (4.578,5.289) and (4.301,5.605) .. (4.043,5.941) .. controls (3.785,6.277) and (3.555,6.633) .. (3.352,7.008) -- cycle
    (9,6.797) .. controls (9.531,6.5) and (9.992,6.113) .. (10.383,5.637) .. controls (10.773,5.16) and (11.062,4.625) .. (11.25,4.031) .. controls (11.5,4.016) and (11.75,4.008) .. (12,4.008) .. controls (12.25,4.008) and (12.5,4.016) .. (12.75,4.031) .. controls (12.938,4.609) and (13.227,5.141) .. (13.617,5.625) .. controls (14.008,6.109) and (14.469,6.5) .. (15,6.797) .. controls (15.531,7.109) and (16.102,7.312) .. (16.711,7.406) .. controls (17.32,7.5) and (17.922,7.484) .. (18.516,7.359) .. controls (18.672,7.562) and (18.812,7.773) .. (18.938,7.992) .. controls (19.062,8.211) and (19.172,8.438) .. (19.266,8.672) .. controls (18.859,9.125) and (18.547,9.641) .. (18.328,10.219) .. controls (18.109,10.797) and (18,11.391) .. (18,12) .. controls (18,12.625) and (18.113,13.227) .. (18.34,13.805) .. controls (18.566,14.383) and (18.875,14.891) .. (19.266,15.328) .. controls (19.172,15.562) and (19.062,15.789) .. (18.938,16.008) .. controls (18.812,16.227) and (18.672,16.438) .. (18.516,16.641) .. controls (17.922,16.516) and (17.32,16.5) .. (16.711,16.594) .. controls (16.102,16.688) and (15.531,16.891) .. (15,17.203) .. controls (14.469,17.5) and (14.008,17.887) .. (13.617,18.363) .. controls (13.227,18.84) and (12.938,19.375) .. (12.75,19.969) .. controls (12.5,19.984) and (12.25,19.992) .. (12,19.992) .. controls (11.75,19.992) and (11.5,19.984) .. (11.25,19.969) .. controls (11.062,19.391) and (10.773,18.859) .. (10.383,18.375) .. controls (9.992,17.891) and (9.531,17.5) .. (9,17.203) .. controls (8.469,16.891) and (7.898,16.688) .. (7.289,16.594) .. controls (6.68,16.5) and (6.078,16.516) .. (5.484,16.641) .. controls (5.328,16.438) and (5.188,16.227) .. (5.062,16.008) .. controls (4.938,15.789) and (4.828,15.562) .. (4.734,15.328) .. controls (5.141,14.875) and (5.453,14.359) .. (5.672,13.781) .. controls (5.891,13.203) and (6,12.609) .. (6,12) .. controls (6,11.375) and (5.887,10.773) .. (5.66,10.195) .. controls (5.434,9.617) and (5.125,9.109) .. (4.734,8.672) .. controls (4.828,8.438) and (4.938,8.211) .. (5.062,7.992) .. controls (5.188,7.773) and (5.328,7.562) .. (5.484,7.359) .. controls (6.078,7.484) and (6.68,7.5) .. (7.289,7.406) .. controls (7.898,7.312) and (8.469,7.109) .. (9,6.797) -- cycle
    (12,9) .. controls (11.609,9) and (11.227,9.074) .. (10.852,9.223) .. controls (10.477,9.371) and (10.148,9.586) .. (9.867,9.867) .. controls (9.586,10.148) and (9.371,10.477) .. (9.223,10.852) .. controls (9.074,11.227) and (9,11.609) .. (9,12) .. controls (9,12.391) and (9.074,12.773) .. (9.223,13.148) .. controls (9.371,13.523) and (9.586,13.852) .. (9.867,14.133) .. controls (10.148,14.414) and (10.477,14.629) .. (10.852,14.777) .. controls (11.227,14.926) and (11.609,15) .. (12,15) .. controls (12.391,15) and (12.773,14.926) .. (13.148,14.777) .. controls (13.523,14.629) and (13.852,14.414) .. (14.133,14.133) .. controls (14.414,13.852) and (14.629,13.523) .. (14.777,13.148) .. controls (14.926,12.773) and (15,12.391) .. (15,12) .. controls (15,11.609) and (14.926,11.227) .. (14.777,10.852) .. controls (14.629,10.477) and (14.414,10.148) .. (14.133,9.867) .. controls (13.852,9.586) and (13.523,9.371) .. (13.148,9.223) .. controls (12.773,9.074) and (12.391,9) .. (12,9) -- cycle
    (12,10.992) .. controls (12.125,10.992) and (12.25,11.02) .. (12.375,11.074) .. controls (12.5,11.129) and (12.609,11.203) .. (12.703,11.297) .. controls (12.797,11.391) and (12.871,11.5) .. (12.926,11.625) .. controls (12.98,11.75) and (13.008,11.875) .. (13.008,12) .. controls (13.008,12.125) and (12.98,12.25) .. (12.926,12.375) .. controls (12.871,12.5) and (12.797,12.609) .. (12.703,12.703) .. controls (12.609,12.797) and (12.5,12.871) .. (12.375,12.926) .. controls (12.25,12.98) and (12.125,13.008) .. (12,13.008) .. controls (11.875,13.008) and (11.75,12.98) .. (11.625,12.926) .. controls (11.5,12.871) and (11.391,12.797) .. (11.297,12.703) .. controls (11.203,12.609) and (11.129,12.5) .. (11.074,12.375) .. controls (11.02,12.25) and (10.992,12.125) .. (10.992,12) .. controls (10.992,11.875) and (11.02,11.75) .. (11.074,11.625) .. controls (11.129,11.5) and (11.203,11.391) .. (11.297,11.297) .. controls (11.391,11.203) and (11.5,11.129) .. (11.625,11.074) .. controls (11.75,11.02) and (11.875,10.992) .. (12,10.992) -- cycle;}

% U+E90C uniE90C – icon-orders (Liste mit Plus); Vorschub 24.000 px; Umriss x 3.00–21.00, y 1.01–19.99 px; 4 Kontur(en)
\newcommand\iconorders{\fill[nonzero rule]
    (18,9) -- (18,6) -- (21,6) -- (21,4.008) -- (18,4.008) -- (18,1.008) -- (16.008,1.008) -- (16.008,4.008) -- (13.008,4.008) -- (13.008,6) -- (16.008,6) -- (16.008,9) -- cycle
    (10.992,6) -- (10.992,4.008) -- (3,4.008) -- (3,6) -- cycle
    (21,13.008) -- (21,10.992) -- (3,10.992) -- (3,13.008) -- (21,13.008) -- cycle
    (21,19.992) -- (21,18) -- (3,18) -- (3,19.992) -- cycle;}

% U+E93A uniE93A – icon-information (Kreis mit i); Vorschub 24.000 px; Umriss x 1.99–22.01, y 1.99–22.01 px; 4 Kontur(en)
\newcommand\iconinfo{\fill[nonzero rule]
    (12,1.992) .. controls (10.625,1.992) and (9.328,2.258) .. (8.109,2.789) .. controls (6.891,3.305) and (5.828,4.016) .. (4.922,4.922) .. controls (4.016,5.828) and (3.305,6.891) .. (2.789,8.109) .. controls (2.258,9.328) and (1.992,10.625) .. (1.992,12) .. controls (1.992,13.375) and (2.258,14.672) .. (2.789,15.891) .. controls (3.305,17.109) and (4.016,18.172) .. (4.922,19.078) .. controls (5.828,19.984) and (6.891,20.695) .. (8.109,21.211) .. controls (9.328,21.742) and (10.625,22.008) .. (12,22.008) .. controls (13.375,22.008) and (14.672,21.742) .. (15.891,21.211) .. controls (17.109,20.695) and (18.172,19.984) .. (19.078,19.078) .. controls (19.984,18.172) and (20.695,17.109) .. (21.211,15.891) .. controls (21.742,14.672) and (22.008,13.375) .. (22.008,12) .. controls (22.008,10.625) and (21.742,9.328) .. (21.211,8.109) .. controls (20.695,6.891) and (19.984,5.828) .. (19.078,4.922) .. controls (18.172,4.016) and (17.109,3.305) .. (15.891,2.789) .. controls (14.672,2.258) and (13.375,1.992) .. (12,1.992) -- cycle
    (12,4.008) .. controls (13.062,4.008) and (14.082,4.211) .. (15.059,4.617) .. controls (16.035,5.023) and (16.898,5.602) .. (17.648,6.352) .. controls (18.398,7.102) and (18.977,7.965) .. (19.383,8.941) .. controls (19.789,9.918) and (19.992,10.938) .. (19.992,12) .. controls (19.992,13.062) and (19.789,14.082) .. (19.383,15.059) .. controls (18.977,16.035) and (18.398,16.898) .. (17.648,17.648) .. controls (16.898,18.398) and (16.035,18.977) .. (15.059,19.383) .. controls (14.082,19.789) and (13.062,19.992) .. (12,19.992) .. controls (10.938,19.992) and (9.918,19.789) .. (8.941,19.383) .. controls (7.965,18.977) and (7.102,18.398) .. (6.352,17.648) .. controls (5.602,16.898) and (5.023,16.035) .. (4.617,15.059) .. controls (4.211,14.082) and (4.008,13.062) .. (4.008,12) .. controls (4.008,10.938) and (4.211,9.918) .. (4.617,8.941) .. controls (5.023,7.965) and (5.602,7.102) .. (6.352,6.352) .. controls (7.102,5.602) and (7.965,5.023) .. (8.941,4.617) .. controls (9.918,4.211) and (10.938,4.008) .. (12,4.008) -- cycle
    (10.992,16.992) -- (13.008,16.992) -- (13.008,15) -- (10.992,15) -- (10.992,16.992) -- cycle
    (10.992,13.008) -- (13.008,13.008) -- (13.008,7.008) -- (10.992,7.008) -- (10.992,13.008) -- cycle;}


\newif\ifmarken \markentrue
\ifdefined\ohnemarken \markenfalse\fi

\newcommand\breite{1484}% 1704 - 70 (Bildspalte) - 150 (Zahnrad, Schließen)
\newcommand\hoehe{145}%   zweite Zeile endet bei 135, + 10 px
\newcommand\links{-70}%   Verschiebung der Tabellenspalten (ID beginnt bei x10)
\newcommand\rand{1474}%   rechter Tabellenrand (\breite - 10)

% Textbreiten in px (vor dem Bild messen; in TikZ ist \nullfont aktiv)
\newlength\tw
\makeatletter
\newcommand\breitepx[2]{\settowidth\tw{\schrift #2}%
  \edef#1{\strip@pt\dimexpr\tw*4/3*7200/7227\relax}}
\newcommand\probe[2]{\breitepx\probewert{#1}%
  \typeout{POSITIONEN-VOLL: '#1' \probewert\space px (gemessen #2)}}
\makeatother
\probe{Active (0)}{55,0}\probe{Pending (0)}{66,1}\probe{History}{40,6}
\probe{ID}{11,9}\probe{Instrument}{60,7}\probe{Type}{28,1}\probe{Amount}{44,1}
\probe{Trade Volume}{78,4}\probe{Open Rate}{59,2}\probe{Open Time}{62,2}
\probe{S/L}{18,8}\probe{T/P}{19,1}\probe{Swap}{31,6}
\probe{Commission}{69,8}\probe{Total Profit}{61,5}\probe{Margin}{39,3}
\probe{All sides}{47,4}\probe{All profits}{53,8}
\breitepx\probewert{Search}\typeout{POSITIONEN-VOLL: 'Search' \probewert\space px (nicht gemessen)}
\breitepx\wn{\halbfett New order}\pgfmathsetmacro\probewert{97.7-\wn-22.3-12}
\typeout{POSITIONEN-VOLL: 'New order' (600) \wn\space px; Abstand Plus-Text \probewert\space px}
\breitepx\wa{Active}\breitepx\wp{Pending}\breitepx\wh{History}
\pgfmathsetmacro\xh{\rand-(\wh+32)}% Reiter rechtsbündig am Tabellenrand
\pgfmathsetmacro\xp{\xh-(\wp+32)}
\pgfmathsetmacro\xa{\xp-(\wa+32)}
\typeout{POSITIONEN-VOLL: Reiter ab x \xa, \xp, \xh\space bis \rand\space px}

% \flaeche{<Zeilenoberkante>}: getönte Zeile, Zellen-Fläche unter dem 1,6-px-Rand;
% außen Radius 3 px (oben innen 3 x 1,4 px)
\newcommand\flaeche[1]{%
  \fill[aktiv,fill opacity=0.1] (\breite/2,#1+1.6) -- (13,#1+1.6)
    arc[start angle=270,end angle=180,x radius=3,y radius=1.4] -- (10,#1+26.6)
    arc[start angle=180,end angle=90,radius=3] -- (\rand-3,#1+29.6)
    arc[start angle=90,end angle=0,radius=3] -- (\rand,#1+3)
    arc[start angle=0,end angle=-90,x radius=3,y radius=1.4] -- cycle;}

% \kopf{<Original-x des span>}{<Text>} / \zelle{<Mitte>}{<y>}{<Farbe>}{<Text>}
\newcommand\kopf[2]{\node[anchor=base west,text=label] at (#1,65.4) {#2};}
\newcommand\zelle[4]{\node[anchor=base,text=#3] at (#1,#2) {#4};}

% \icon{<x>}{<y Oberkante>}{<Größe px>}{<Farbe>}{<Icon>}: Box Größe x Größe
% (icons.tikz: 24 px, Ursprung unten links, y nach oben)
\newcommand\icon[5]{%
  \begin{scope}[shift={(#1,#2+#3)},x=#3/32*1bp,y=#3/32*1bp,color=#4]#5\end{scope}}

% \marke{<Original-x des Spaltennamens>}{<Nummer>}: Kreis 17 px, Mitte 3 px + Radius davor
\newcommand\marke[2]{%
  \fill[white] (#1-11.5,61) circle[radius=8.5];
  \node[anchor=base,text=markenziffer,font=\fontsize{6.954bp}{8bp}\halbfett]
    at (#1-11.5,61+3.091) {#2};}% 12 px * 17/22

\begin{document}%
\begin{tikzpicture}[x=\px,y={(0,-\px)},font=\schrift,
    every node/.style={inner sep=0,outer sep=0}]
  \path[use as bounding box] (0,0) rectangle (\breite,\hoehe);
  % Block und Zeilen
  \fill[block,rounded corners=7\px] (0,0) rectangle (\breite,\hoehe);
  \flaeche{75.8}\flaeche{105.4}
  % Leiste oben links (Lage wie gemessen)
  \draw[gewinn,line width=0.8\px,rounded corners=4.6\px]
    (10.4,10.6) rectangle (109.6,39.4);% Rand 0,8 px innen, Radius 5 px außen
  \ifdefined\iconplus \icon{22.3}{17.8}{12}{white}{\iconplus}\fi
  \node[anchor=base east,text=white,font=\schrift\halbfett] at (97.7,29.0) {New order};
  \fill[reiter,rounded corners=4\px] (120,10) rectangle (265,40);
  \node[anchor=base west,text=suchtext] at (130.8,29.2) {Search};
  \ifdefined\iconsearch \icon{238}{17.5}{15}{label}{\iconsearch}\fi
  \fill[reiter,rounded corners=4\px] (275,10) rectangle (420,40);
  \node[anchor=base west,text=white] at (285.8,29.4) {All sides};
  \fill[reiter,rounded corners=4\px] (430,10) rectangle (575,40);
  \node[anchor=base west,text=white] at (440.8,29.4) {All profits};
  % Reiter (ohne Zähler), bündig mit dem rechten Tabellenrand
  \fill[reiter,rounded corners=3\px] (\xa,10) rectangle (\rand,40);
  \fill[aktiv] (\xp,10) {[rounded corners=3\px] -- (\xa,10) -- (\xa,40)} -- (\xp,40) -- cycle;
  \node[anchor=base,text=white] at ({\xa+(\wa+32)/2},29.2) {Active};
  \node[anchor=base,text=white] at ({\xp+(\wp+32)/2},29.2) {Pending};
  \node[anchor=base,text=white] at ({\xh+(\wh+32)/2},29.2) {History};
  % Tabelle (Original-x, um die Bildspalte nach links verschoben)
  \begin{scope}[xshift=\links\px]
    \kopf{102.5}{ID}\kopf{177}{Instrument}\kopf{298}{Type}\kopf{383.8}{Amount}
    \kopf{493.5}{Trade Volume}\kopf{640}{Open Rate}\kopf{787.5}{Open Time}
    \kopf{941}{S/L}\kopf{1030.9}{T/P}\kopf{1113.6}{Swap}\kopf{1207.4}{Commission}
    \kopf{1344.2}{Total Profit}\kopf{1467.9}{Margin}
    \foreach \y/\farbe in {95.0/gewinn,124.6/verlust} {
      \zelle{115.9}{\y}{white}{Nummer}
      \zelle{209.4}{\y}{white}{Instrument}% Mitte 214,9 - 5,5 (Pfeil-Bild weggelassen)
      \zelle{319.6}{\y}{white}{Buy}
      \zelle{413.35}{\y}{white}{Menge}
      \zelle{540.25}{\y}{white}{Wert}
      \zelle{677.1}{\y}{white}{Kurs}
      \zelle{826.1}{\y}{white}{Zeitpunkt}
      \zelle{957.9}{\y}{white}{—}
      \zelle{1047.9}{\y}{white}{—}
      \zelle{1136.9}{\y}{white}{Betrag}
      \zelle{1249.85}{\y}{white}{—}
      \zelle{1382.5}{\y}{\farbe}{Ergebnis}
      \zelle{1495}{\y}{white}{Betrag}}
    \ifmarken
      \marke{1113.6}{1}
      \marke{1344.2}{2}
      \marke{1467.9}{3}
    \fi
  \end{scope}
\end{tikzpicture}%
\end{document}
